
% ----------------------------------------------------------------------
%                   LATEX TEMPLATE FOR PhD THESIS
% ----------------------------------------------------------------------

% based on Harish Bhanderi's PhD/MPhil template, then Uni Cambridge
% http://www-h.eng.cam.ac.uk/help/tpl/textprocessing/ThesisStyle/
% corrected and extended in 2007 by Jakob Suckale, then MPI-CBG PhD programme
% and made available through OpenWetWare.org - the free biology wiki


%: Style file for Latex
% Most style definitions are in the external file PhDthesisPSnPDF.
% In this template package, it can be found in ./Latex/Classes/
\documentclass[twoside,11pt]{Latex/Classes/PhDthesisPSnPDF}


%: Macro file for Latex
% Macros help you summarise frequently repeated Latex commands.
% Here, they are placed in an external file /Latex/Macros/MacroFile1.tex
% An macro that you may use frequently is the figuremacro (see introduction.tex)
% This file contains macros that can be called up from connected TeX files
% It helps to summarise repeated code, e.g. figure insertion (see below).

% insert a centered figure with caption and description
% parameters 1:filename, 2:title, 3:description and label
\newcommand{\figuremacro}[3]{
	\begin{figure}[htbp]
		\centering
		\includegraphics[width=1\textwidth]{#1}
		\caption[#2]{\textbf{#2} - #3}
		\label{#1}
	\end{figure}
}

% insert a centered figure with caption and description AND WIDTH
% parameters 1:filename, 2:title, 3:description and label, 4: textwidth
% textwidth 1 means as text, 0.5 means half the width of the text
\newcommand{\figuremacroW}[4]{
	\begin{figure}[htbp]
		\centering
		\includegraphics[width=#4\textwidth]{#1}
		\caption[#2]{\textbf{#2} - #3}
		\label{#1}
	\end{figure}
}

% inserts a figure with wrapped around text; only suitable for NARROW figs
% o is for outside on a double paged document; others: l, r, i(inside)
% text and figure will each be half of the document width
% note: long captions often crash with adjacent content; take care
% in general: above 2 macro produce more reliable layout
\newcommand{\figuremacroN}[3]{
	\begin{wrapfigure}{o}{0.5\textwidth}
		\centering
		\includegraphics[width=0.48\textwidth]{#1}
		\caption[#2]{{\small\textbf{#2} - #3}}
		\label{#1}
	\end{wrapfigure}
}

% predefined commands by Harish
\newcommand{\PdfPsText}[2]{
  \ifpdf
     #1
  \else
     #2
  \fi
}

\newcommand{\IncludeGraphicsH}[3]{
  \PdfPsText{\includegraphics[height=#2]{#1}}{\includegraphics[bb = #3, height=#2]{#1}}
}

\newcommand{\IncludeGraphicsW}[3]{
  \PdfPsText{\includegraphics[width=#2]{#1}}{\includegraphics[bb = #3, width=#2]{#1}}
}

\newcommand{\InsertFig}[3]{
  \begin{figure}[!htbp]
    \begin{center}
      \leavevmode
      #1
      \caption{#2}
      \label{#3}
    \end{center}
  \end{figure}
}


%%% Local Variables: 
%%% mode: latex
%%% TeX-master: "~/Documents/LaTeX/CUEDThesisPSnPDF/thesis"
%%% End: 


 \usepackage{amsmath}

%: ----------------------------------------------------------------------
%:                  TITLE PAGE: name, degree,..
% ----------------------------------------------------------------------
% below is to generate the title page with crest and author name

%if output to PDF then put the following in PDF header
\ifpdf  
    \pdfinfo { /Title  (PhD and MPhil Thesis Classes)
               /Creator (TeX)
               /Producer (pdfTeX)
               /Author (YourName your@email.net)
               /CreationDate (D:YYYYMMDDhhmmss)  %format D:YYYYMMDDhhmmss
               /ModDate (D:YYYYMMDDhhmm)
               /Subject (xyz)
               /Keywords (add, your, keywords, here) }
    \pdfcatalog { /PageMode (/UseOutlines)
                  /OpenAction (fitbh)  }
\fi


\title{Instrumentation}



% ----------------------------------------------------------------------
% The section below defines www links/email for author and institutions
% They will appear on the title page of the PDF and can be clicked
\ifpdf
  \author{\href{mailto:20261142@student.uwa.edu.au}{Beau Trepp}}
%  \cityofbirth{born in XYZ} % uncomment this if your university requires this
%  % If city of birth is required, also uncomment 2 sections in PhDthesisPSnPDF
%  % Just search for the "city" and you'll find them.
  \collegeordept{\href{http://ee.uwa.edu.au}{Electrical, Electronic and Computer Engineering}}
  \university{\href{http://www.uwa.edu.au}{University Of Western Australia}}

  % The crest is a graphics file of the logo of your research institution.
  % Place it in ./0_frontmatter/figures and specify the width
  \crest{\includegraphics[width=4cm]{UWA_crest}}
  
% If you are not creating a PDF then use the following. The default is PDF.
\else
  \author{Beau Trepp}
%  \cityofbirth{born in XYZ}
  \collegeordept{Electrical, Electronic and Computer Engineering}
  \university{University Of Western Australia}
  \crest{\includegraphics[width=4cm]{UWA_crest}}
\fi

%\renewcommand{\submittedtext}{change the default text here if needed}
\degree{Bachelor of Computer Science/ Bachelor of Electronic Engineering}
\degreedate{2011 November}


% ----------------------------------------------------------------------
       
% turn of those nasty overfull and underfull hboxes
\hbadness=10000
\hfuzz=50pt



%: --------------------------------------------------------------
%:                  FRONT MATTER: dedications, abstract,..
% --------------------------------------------------------------

\begin{document}

%\language{english}

% sets line spacing
\renewcommand\baselinestretch{1.2}
\baselineskip=18pt plus1pt


%: ----------------------- generate cover page ------------------------

\maketitle  % command to print the title page with above variables


%: ----------------------- cover page back side ------------------------
% Your research institution may require reviewer names, etc.
% This cover back side is required by Dresden Med Fac; uncomment if needed.

%\newpage
%\vspace{10mm}
%1. Reviewer: Name
%
%\vspace{10mm}
%2. Reviewer: 
%
%\vspace{20mm}
%Day of the defense:
%
%\vspace{20mm}
%\hspace{70mm}Signature from head of PhD committee:



%: ----------------------- abstract ------------------------

% Your institution may have specific regulations if you need an abstract and where it is to be placed in the document. The default here is just after title.


% Thesis Abstract -----------------------------------------------------


%\begin{abstractslong}    %uncommenting this line, gives a different abstract heading
\begin{abstracts}        %this creates the heading for the abstract page

The topic of electric vehicles is becoming increasingly popular due to rising fuel costs and growing concern over emissions. Despite this attention, most electric vehicles have little or no telemetry systems, making many aspects of there operation and efficiency a mystery. 

The aim of this project was to develop an extend-able system in order to capture various data-points that can be available in a vehicle, as well as an interface to display this data inside the car. The design developed differs from traditional embedded systems by being completely modular. It uses existing network protocols to allow the system to be distributed between various smaller embedded components. This will enable it to be easily extended, should the need for more data-points arise, and allows the use of many smaller systems to be implemented incrementally, rather than one expensive monolithic design.

In order to help facilitate robustness and code re-use, a windowing toolkit was also developed. This provides a common platform for user interaction, and defines pre-built components in order to speed up the design and implementations of the user interface.

By exposing and recording more data, deeper analysis can be done on the efficiency of the car, and help justify different technological improvements to the vehicle. The higher granularity of data acquired can also be used to analyze the economy of the vehicle in different conditions.

\end{abstracts}
%\end{abstractlongs}


% ---------------------------------------------------------------------- 


% The original template provides and abstractseparate environment, if your institution requires them to be separate. I think it's easier to print the abstract from the complete thesis by restricting printing to the relevant page.
% \begin{abstractseparate}
%   
% Thesis Abstract -----------------------------------------------------


%\begin{abstractslong}    %uncommenting this line, gives a different abstract heading
\begin{abstracts}        %this creates the heading for the abstract page

The topic of electric vehicles is becoming increasingly popular due to rising fuel costs and growing concern over emissions. Despite this attention, most electric vehicles have little or no telemetry systems, making many aspects of there operation and efficiency a mystery. 

The aim of this project was to develop an extend-able system in order to capture various data-points that can be available in a vehicle, as well as an interface to display this data inside the car. The design developed differs from traditional embedded systems by being completely modular. It uses existing network protocols to allow the system to be distributed between various smaller embedded components. This will enable it to be easily extended, should the need for more data-points arise, and allows the use of many smaller systems to be implemented incrementally, rather than one expensive monolithic design.

In order to help facilitate robustness and code re-use, a windowing toolkit was also developed. This provides a common platform for user interaction, and defines pre-built components in order to speed up the design and implementations of the user interface.

By exposing and recording more data, deeper analysis can be done on the efficiency of the car, and help justify different technological improvements to the vehicle. The higher granularity of data acquired can also be used to analyze the economy of the vehicle in different conditions.

\end{abstracts}
%\end{abstractlongs}


% ---------------------------------------------------------------------- 

% \end{abstractseparate}


%: ----------------------- tie in front matter ------------------------

\frontmatter
% Thesis Dedictation ---------------------------------------------------

\begin{dedication} %this creates the heading for the dedication page

%To ...

\begin{flushright}
 

Beau Trepp 

28 Waterfall Rd

Forrestfield WA 6058 
\end{flushright}

\begin{flushleft}
The Dean
Faculty of Engineering Computing and Mathematics \newline
The University of Western Australia \newline
35 Stirling Highway \newline
CRAWLEY WA 6009 
\end{flushleft}

\begin{flushleft}
Dear Sir, \newline
  \newline
I submit to you this dissertation entitled ― "Design of an embedded data acquisition and display system using a modular approach" in partial fulfillment of the requirement of the award of Bachelor of Engineering/Bachelor of Computer Science.
 \newline \newline \newline
Yours Faithfully, \newline
\newline \newline
Beau Trepp \newline
20261142
\end{flushleft}
\end{dedication}

% ----------------------------------------------------------------------
% Thesis Acknowledgements ------------------------------------------------


%\begin{acknowledgementslong} %uncommenting this line, gives a different acknowledgements heading
\begin{acknowledgements}      %this creates the heading for the acknowlegments

During work on this project, I received help from many people whom I would like to thank.

The members of the gumstix and ZeroMQ mailing lists who helped me with issues regarding compilation of both code-bases.

Ian Hooper who helped with work on the vehicle, and provided knowledge into the construction of the vehicle.

Jonathan Oakley who showed helped me with understanding the safety issues regarding driving the car, and for showing me the routes used to performance test it.

The maintainers of the open icon library, for the images used in the project.

To my supervisor Prof Thomas Bräunl for the opportunity to work on the project and vehicle and for guidance during the project.

Kevin Vinsen, for looking after the REV project while Thomas was on sabbatical.

Finally I would like to thank all the members of the REV Team, for providing assistance and support throughout the project.

\end{acknowledgements}
%\end{acknowledgmentslong}

% ------------------------------------------------------------------------





%: ----------------------- contents ------------------------

\setcounter{secnumdepth}{3} % organisational level that receives a numbers
\setcounter{tocdepth}{3}    % print table of contents for level 3
\tableofcontents            % print the table of contents
% levels are: 0 - chapter, 1 - section, 2 - subsection, 3 - subsection


%: ----------------------- list of figures/tables ------------------------

\listoffigures	% print list of figures

\listoftables  % print list of tables


%: ----------------------- glossary ------------------------

% Tie in external source file for definitions: /0_frontmatter/glossary.tex
% Glossary entries can also be defined in the main text. See glossary.tex
% this file is called up by thesis.tex
% content in this file will be fed into the main document

% Glossary entries are defined with the command \nomenclature{1}{2}
% 1 = Entry name, e.g. abbreviation; 2 = Explanation
% You can place all explanations in this separate file or declare them in the middle of the text. Either way they will be collected in the glossary.

% required to print nomenclature name to page header
\markboth{\MakeUppercase{\nomname}}{\MakeUppercase{\nomname}}


% ----------------------- contents from here ------------------------

% chemicals
\nomenclature{DAPI}{4',6-diamidino-2-phenylindole; a fluorescent stain that binds strongly to DNA and serves to marks the nucleus in fluorescence microscopy} 
\nomenclature{DEPC}{diethyl-pyro-carbonate; used to remove RNA-degrading enzymes (RNAases) from water and laboratory utensils}
\nomenclature{DMSO}{dimethyl sulfoxide; organic solvent, readily passes through skin, cryoprotectant in cell culture}
\nomenclature{EDTA}{Ethylene-diamine-tetraacetic acid; a chelating (two-pronged) molecule used to sequester most divalent (or trivalent) metal ions, such as calcium (Ca$^{2+}$) and magnesium (Mg$^{2+}$), copper (Cu$^{2+}$), or iron (Fe$^{2+}$ / Fe$^{3+}$)}



 

\begin{multicols}{2} % \begin{multicols}{#columns}[header text][space]
\begin{footnotesize} % scriptsize(7) < footnotesize(8) < small (9) < normal (10)

\printnomenclature[1.5cm] % [] = distance between entry and description
\label{nom} % target name for links to glossary

\end{footnotesize}
\end{multicols}



%: --------------------------------------------------------------
%:                  MAIN DOCUMENT SECTION
% --------------------------------------------------------------

% the main text starts here with the introduction, 1st chapter,...
\mainmatter

\renewcommand{\chaptername}{} % uncomment to print only "1" not "Chapter 1"


%: ----------------------- subdocuments ------------------------

% Parts of the thesis are included below. Rename the files as required.
% But take care that the paths match. You can also change the order of appearance by moving the include commands.


% this file is called up by thesis.tex
% content in this file will be fed into the main document

%: ----------------------- introduction file header -----------------------
\chapter{Introduction}

% the code below specifies where the figures are stored
\ifpdf
    \graphicspath{{1_introduction/figures/PNG/}{1_introduction/figures/PDF/}{1_introduction/figures/}}
\else
    \graphicspath{{1_introduction/figures/EPS/}{1_introduction/figures/}}
\fi

% ----------------------------------------------------------------------
%: ----------------------- introduction content ----------------------- 
% ----------------------------------------------------------------------



%: ----------------------- HELP: latex document organisation
% the commands below help you to subdivide and organise your thesis
%    \chapter{}       = level 1, top level
%    \section{}       = level 2
%    \subsection{}    = level 3
%    \subsubsection{} = level 4
% note that everything after the percentage sign is hidden from output



\section{put section name here} % section headings are printed smaller than chapter names
% intro
Write your text without any further commands, like this:.... Any organised system requires energy, be it a machine of some kind or a live organism. Energy is needed to win the uphill battle against entropy and pull together lifeless molecules to be able to do something in this world, like complete a PhD.



\subsection{Name your subsection} % subsection headings are again smaller than section names
% lead
Different organised systems have different energy currencies. The machines that enable us to do science like sizzling electricity but at a controlled voltage. Earth's living beings are no different, except that they have developed another preference. They thrive on various chemicals. 

% dextran, starch, glycogen
Most organisms use polymers of glucose units for energy storage and differ only slightly in the way they link together monomers to sometimes gigantic macromolecules. Dextran of bacteria is made from long chains of $\alpha$-1,6-linked glucose units. 

%: ----------------------- HELP: special characters
% above you can see how special characters are coded; e.g. $\alpha$
% below are the most frequently used codes:
%$\alpha$  $\beta$  $\gamma$  $\delta$

%$^{chars to be superscripted}$  OR $^x$ (for a single character)
%$_{chars to be suberscripted}$  OR $_x$

%>  $>$  greater,  <  $<$  less
%≥  $\ge$  greater than or equal, ≤  $\ge$  lesser than or equal
%~  $\sim$  similar to

%$^{\circ}$C   ° as in degree C
%±  \pm     plus/minus sign

%$\AA$     produces  Å (Angstrom)




% dextran, starch, glycogen continued
Starch of plants and glycogen of animals consists of $\alpha$-1,4-glycosidic glucose polymers \cite{lastname07}. See figure \ref{largepotato} for a comparison of glucose polymer structure and chemistry. 

Two references can be placed separated by a comma \cite{lastname07,name06}.

%: ----------------------- HELP: references
% References can be links to figures, tables, sections, or references.
% For figures, tables, and text you define the target of the link with \label{XYZ}. Then you call cross-link with the command \ref{XYZ}, as above
% Citations are bound in a very similar way with \cite{XYZ}. You store your references in a BibTex file with a programme like BibDesk.





\figuremacro{largepotato}{A common glucose polymers}{The figure shows starch granules in potato cells, taken from \href{http://molecularexpressions.com/micro/gallery/burgersnfries/burgersnfries4.html}{Molecular Expressions}.}

%: ----------------------- HELP: adding figures with macros
% This template provides a very convenient way to add figures with minimal code.
% \figuremacro{1}{2}{3}{4} calls up a series of commands formating your image.
% 1 = name of the file without extension; PNG, JPEG is ok; GIF doesn't work
% 2 = title of the figure AND the name of the label for cross-linking
% 3 = caption text for the figure

%: ----------------------- HELP: www links
% You can also see above how, www links are placed
% \href{http://www.something.net}{link text}

\figuremacroW{largepotato}{Title}{Caption}{0.8}
% variation of the above macro with a width setting
% \figuremacroW{1}{2}{3}{4}
% 1-3 as above
% 4 = size relative to text width which is 1; use this to reduce figures




Insulin stimulates the following processes:

\begin{itemize}
\item muscle and fat cells remove glucose from the blood,
\item cells breakdown glucose via glycolysis and the citrate cycle, storing its energy in the form of ATP,
\item liver and muscle store glucose as glycogen as a short-term energy reserve,
\item adipose tissue stores glucose as fat for long-term energy reserve, and
\item cells use glucose for protein synthesis.
\end{itemize}

%: ----------------------- HELP: lists
% This is how you generate lists in LaTeX.
% If you replace {itemize} by {enumerate} you get a numbered list.


 


%: ----------------------- HELP: tables
% Directly coding tables in latex is tiresome. See below.
% I would recommend using a converter macro that allows you to make the table in Excel and convert them into latex code which you can then paste into your doc.
% This is the link: http://www.softpedia.com/get/Office-tools/Other-Office-Tools/Excel2Latex.shtml
% It's a Excel template file containing a macro for the conversion.

\begin{table}[htdp]
\centering
\begin{tabular}{ccc} % ccc means 3 columns, all centered; alternatives are l, r

{\bf Gene} & {\bf GeneID} & {\bf Length} \\ 
% & denotes the end of a cell/column, \\ changes to next table row
\hline % draws a line under the column headers

human latexin & 1234 & 14.9 kbps \\
mouse latexin & 2345 & 10.1 kbps \\
rat latexin   & 3456 & 9.6 kbps \\
% Watch out. Every line must have 3 columns = 2x &. 
% Otherwise you will get an error.

\end{tabular}
\caption[title of table]{\textbf{title of table} - Overview of latexin genes.}
% You only need to write the title twice if you don't want it to appear in bold in the list of tables.
\label{latexin_genes} % label for cross-links with \ref{latexin_genes}
\end{table}



% There you go. You already know the most important things.


% ----------------------------------------------------------------------



	% background information
% this file is called up by thesis.tex
% content in this file will be fed into the main document

\chapter{Aims of the project} % top level followed by section, subsection


% ----------------------- contents from here ------------------------

\section{Final Aim}

The Ultimate goal of the project is to developed a distributed system for use in an automotive environment. This will culminate into a completed system, provided data logging functionality and a user interface to view the live data.

\section{Goals}

The goals this project aimed to achieve were set out as follows.

\begin{enumerate}
\item  Investigate the use messaging protocols inside a minimal embedded systems
\item  Develop GPS capability
\item  Develop Battery monitor capabilty
\item  Integrate this data into a user display
\item  Log this data to be reviewed later
\end{enumerate}

During development, the author also kept the following goals in-line, to make the system easier to develop for.  This was done to give the software a better chance of being developed further by future projects on this topic.


\begin{enumerate}
\item  Allow the system to be easily extend-able
\item  Attempt to minimize future exposure to concurrency
\item  Use simpler individual designs, to make it easier for engineers (as opposed to programmers) to modify
\item Develop standards for inter-module communication
\item Develop common UI elements to facilitate ease-of-layout and interface design.
\item Develop the system in such a way that it is portable to other devices. The hardware used is slightly older, and the system will be developed in a way that it can be ported across to new hardware.
\end{enumerate}


% ---------------------------------------------------------------------------
% ----------------------- end of thesis sub-document ------------------------
% ---------------------------------------------------------------------------						% aims of the project

% this file is called up by thesis.tex
% content in this file will be fed into the main document

%: ----------------------- name of chapter  -------------------------
\chapter{Literature Review} % top level followed by section, subsection


%: ----------------------- paths to graphics ------------------------

% change according to folder and file names
\ifpdf
    \graphicspath{{X/figures/PNG/}{X/figures/PDF/}{X/figures/}}
\else
    \graphicspath{{X/figures/EPS/}{X/figures/}}
\fi

%: ----------------------- contents from here ------------------------


\section{Renewable Energy Vehicle Instrumentation:
Graphical User Interface and Black Box}

In 2009 Daksh Varma completed a project entitled \emph{'Renewable Energy Vehicle Instrumentation:
Graphical User Interface and Black Box'} \cite{thesis_varma}. This project developed systems for the Hyundai Getz featured in this project, and another vehicle, a Lotus Elise. This project focused on developing user interfaces and logging functionality for both vehicles.

The system for the Getz was implemented on the current hardware installed in the vehicle. Varma implemented everything using C, rather than experimenting with more complicated languages. This code was very much tied to the current hardware configuration of the car, and is not functional with the unknown hardware changes that have occurred to the car since 2009. This highlights an interesting point, even though Varma had created a functional system, changes that happened after his system was deployed had large negative effects. Changes that were not under Varma's control, caused the system to crash and hang. This decreased the acceptance and view of the system.

Varma used the FPGA present on the board to interface with the digital outputs available. This implementation would produce the fastest speed, but has some problems. Due to differences in the hardware revisions of the eyebot, several eyebots require different drivers for the FPGA. This means any code developed would require the correct drivers, which can be difficult to maintain. It also greatly ties the program to the eyebot platform, which limits future expansion to other machines.

For threading Varma used the timers that were implemented in the RobiOS libraries. These would cause functions to be called at periodic intervals. Varma discussed the use of locks in order to protect shared variables. This of course did make the system complicated. Locking variables in thread's, if not implemented correctly, has a chance of causing deadlocks. This can cause the program to never respond to user input. The use of timers also restricts the ability of implementing the code on platforms that do not run RobiOS.

When developing code for the Lotus, Varma separated front-end and back-end functionality. This separation was a good idea, as it helped provide robustness. The use of two programming languages allowed Varma to leverage the strengths of each individually.

When implementing the maps functionality in the Lotus, Varma created the whole of Perth as one giant image. This would make the code simpler to maintain, as there is no complicated image swapping that needs to be performed. However this greatly limits the usability of the map system developed. The giant map image would need to be replaced manually for driving the car in different locations. It would also consume more ram than is necessary, as the whole image is always in memory. The map system develop also has no support for zoom, leaving the display stuck at the one size.

Varma's system was functional and performed its tasks well. It suffers from a very rigid design, that is not very expandable. While a good technique of separating the user interface and the backend functionality was implemented for one vehicle, it was not done for the other. There is no indication why this was the case. Varma's system implemented functionality not seen on other cars, and improved the experience for the driver.

\section{Development of a User Interface
for Electric Cars}

In 2010 Thomas Walter undertook a project names \emph{'Development of a User Interface
for Electric Cars'}. This project was concerned with developing a user interface for the Lotus Elise.

Walter implemented the interface by using a windowing toolkit named Qt. This toolkit is available for all major platforms, such as the windows operating system the project was developed on. Qt provides a large array of elements that the designer can use to lay-out the interface. The use of this framework would make development easier, and the final product more stable.

Like Varma in 2009, Walter continued the ideal of separating the user interface from the back-end of the program. This helped diagnose problems with the system, and provided increased flexibility in development.

Walter heavily relied on the use of an Object Orientated language in order to implement his system. This abstraction made it easier to identify what each classes functionality was, and makes the overall system design easier to understand. This also tied in well with the use of Qt for the user interface, as Qt is written in an Object Orientated Language.

Network logging was an interesting feature added in Walters project. This logger enable the system to log data via an Internet enabled server. This logger functionality takes into account the unreliable nature of the connection used, and stores unsent messages in a queue if communication is lost.

Walter provides a solution to the problems with the map system used in Varma's project. Rather than storing the map in one giant file, Walter uses a technique called tiling. This technique breaks the map up into many smaller images. By knowing the current location, the relevant images can be loaded and displayed. This reduces memory consumption, as only the active tiles are displayed, and allows the system to expand to other locations.

Walter's design utilized abstractions and data structures in order to provide a more expandable system. This enabled the system to adapt to changes in location, and problems with the network. These two features are highly desirable, and would be welcome additions to the system installed in the Getz.




% ---------------------------------------------------------------------------
%: ----------------------- end of thesis sub-document ------------------------
% ---------------------------------------------------------------------------

			
% this file is called up by thesis.tex
% content in this file will be fed into the main document

%: ----------------------- name of chapter  -------------------------
\chapter{System Design} % top level followed by section, subsection


%: ----------------------- paths to graphics ------------------------

% change according to folder and file names
\ifpdf
    \graphicspath{{4/figures/PNG/}{4/figures/PDF/}{4/figures/}}
\else
    \graphicspath{{4/figures/EPS/}{4/figures/}}
\fi

%: ----------------------- contents from here ------------------------



%: ----------------------- contents from here ------------------------

\section{GPS Module}

\subsection{Hardware}

A vital part of the data logging and user-interface of the software is finding the cars current location. This is done by the use of a “off-the self” GPS unit. Currently the system is using a Qstar usb equipped GPS receiver. This receiver operates at a rate of 10hz , though it can be set to operate at a slower frequency of 1hz. For the purposes of recording positional data, along with estimating the vehicles current speed, the unit should run as fast as possible. The extra precision is useful for the data-logging aspect, with no negative effects on the user-display aspect.  

The GPS device can be enumerated as a standard serial port. This is beneficial as it can be used on any device that has the correct drivers and a available usb port. As it appears as a normal serial port, it can be queried using standard system routines. This allows the program that reads the device to operate any custom knowledge of the device it is connected to, aside from the serial parameters to make the connection. 

\subsection{ Drivers}
\label{sec:gpsdrivers}

While the Eyebot M6 has hardware usb support, it was not immediately compatible with the GPS sensor. Various versions of usb-serial drivers where tried each with their own problems. The main cause of this difficultly was the out-dated Linux kernel being run in the system. This was kernel version number 2.6.17 and was released in 2006, which is 5 years old as of writing \cite{osnews}. This was a major cause of incompatibilities, as the GPS receiver was manufactured a significant time after this kernel was written. The drivers had no clue as to what the usb product keys were, nor the specific quirks that the devices may have had.

The first driver attempted was the generic usb-serial driver, included as a kernel module in 2.6.17. This drivers success would mean that the sensor and program could be easily installed in just about any machine running Linux. The driver would have matured after the 2.6.17 kernel, and newer kernels would have support by default. This is beneficial to the system as it would require the least amount of configuration and setup if the GPS program was set to run in a different machine.

Sadly this driver did not perform correctly with the GPS device. While the driver was able to be loaded into the kernel without any errors, it caused problems when trying to associate with the GPS. The device uses extra features, which were unsupported by the generic driver. This caused strange symptoms in the operating system. The main symptom of an incorrect driver was the generation of the /dev/ttyUSB0 device. This availability of this device implies that a tty is available to read/write from. Due to the incompatibility of the driver, this serial port would never report any bytes to be read, which is why it is unsuitable for use with this device. Customizing the generic driver to support this device would be unfeasible because it is unlikely that newer versions of this driver would support the device. This leads to the situation that if the GPS program is ported to a different machine, a custom version of the generic usb serial drivers would have to be ported as well.

As the GPS device did not work with the generic serial drivers, alternative drivers were investigated in order to support this device. Experiments indicated that this device was automatically detected and loaded in a newer kernel. This functioned correctly and was able to communicate with the GPS device at the full rate of 10hz. The driver used by this kernel was called cdc-acm. Further investigation showed that this driver could be included as a kernel module for the gumstix platform.

This driver was not immediately compatible with the device. This was because the device was manufactured after the kernel . As such the driver did not recognize the manufacturer ID and product ID of the GPS. The driver was then modified to include this information and re-deployed to the eye-bot. This was successful in creating the virtual serial device inside /dev, and also in allowing data to be read from this device.

While the cdc-acm driver was able to be loaded and functioned, it still contained errors. If the system was under intense CPU load, the program may not run quick enough to remove all the data from the serial port buffer. This would cause the operating system to throttle the port. Examination of the driver source code reveals that new information is dropped while this driver is throttled. This is acceptable behavior in this instance, however this driver had a race condition. If the TTY was throttled under certain conditions, it would be unable to unthrottle later on. This cause the TTY to drain its buffer and never accept any new data from the GPS even if its buffer was empty. The driver was further modified to include spin-locks, a primitive kernel locking technique, in order to prevent this situation. The driver is now able to run for extended periods of time without locking up, enabling a reliable GPS reporting mechanism to be developed. 

\subsection {Design}

Development of the GPS reporting component was done in C. This was chosen as it is a relatively low level language, with wide support. It is simpler to understand than more complicated object oriented style languages. This makes it a good choice for the GPS reporting mechanism, as it only has to do one task. In order to ensure that the code can be easily modified by future programmers, the structure of this program is simple. It runs in only one-thread, aside from the back ground ZeroMQ threads, and thus requires no concurrency management.

\subsubsection{Initial Design}

This first iteration of this code used blocking ports and read a single byte at a time. This was the style used to match existing examples \cite{thesis_varma}. This was a functional design however it did lead to some problems. One problem with this approach was excessive throttling. The process would be woken up every time one character could be read, and would only remove one character from the buffer, even if there were hundreds waiting. This is a bad situation, as the process will continually be awoken to do trivial work. This steals cpu-time from another process, and caused the program to appear “sluggish”. This design was also in-sufficient as blocked the process while attempting to read from the port. This would cause the program to appear to hang if no data was available. It made it difficult to diagnose errors with this program.

\subsubsection{Final Design}

The design was refined in order to support bulk-reads and non-blocking operation. This fixes the two problems with the first approach. Rather than reading one character at a time, the program now reads as many as possible and stores the information into it's own circular buffer. This allows the serial port to be purged as quickly as possible. It also has the added feature of allowing the program to decide how to discard messages in the case that it cannot keep up with the GPS. Figure \ref{gpsflowchart} shows the program flow of the final design.

\figuremacro{gpsflowchart}{Flow chart of the GPS daemon}{}


\subsubsection{Network Protocol}

Table \ref{tab:gpsprotocol} shows the protocol for the GPS message when transmitted over the network. The protocol uses a binary format instead of an ASCII based one. This reduces the space/data transmitted over the link, which helps reduce cost and improve speed. The motivation for ASCII based protocols is that control characters can be used to help synchronize the data. As this design uses ZeroMQ in order to manage the flow of data, such control characters are unnecessary. All values are transmitted in network order, this is big-endian order so that the most significant byte is transmitted first.

\begin{table}
    \begin{tabular}{*{6}{|l}|}
        \hline
        47 &  50 & 53 & 0 & 0  & 01    \\ \hline \hline
        \multicolumn{5}{|c|}{ Header} & \multicolumn {1}{|c|}{Data Payload} \\ \hline
	\multicolumn{3}{|c|}{Filter} & Minor Filter & Revision Number & Fix \\
        \hline
    \end{tabular}
\end{table}

\begin{table}
    \begin{tabular}{*{6}{|l}|}
        \hline
        4E1FE968 & 00000000 & C1FFD54D & 42E7A1F1 & 43B40000 & 41200000 \\ \hline \hline
         \multicolumn {6}{|c|}{Data Payload} \\ \hline
	Time(s) & Time(us) & Latitude & Longitude & Bearing & Speed \\
        \hline
    \end{tabular}
	\caption{Network protocol for GPS daemon}
	\label{tab:gpsprotocol}
\end{table}



% ---------------------------------------------------------------------------
%: ----------------------- end of thesis sub-document ------------------------
% ---------------------------------------------------------------------------



%: ----------------------- contents from here ------------------------

\section{TBS Module}

\subsection{Hardware}

The most important external device used in the user interface and data-logging aspect of the software is that of the battery monitoring module. The car has 45 Lithium Ion batteries installed, and it is useful to monitor the charge, current and voltage of the battery cells at all times. The system that the monitoring software runs on is not a highly reliable embedded system. It requires a few minutes to start up, and consumes too much power to leave running all the time. As such a different device is used to track the health and charge of the batteries. This device is a e-xpert pro battery monitor manufactured by TBS electronics. This is a commerical unit which increases the reliability of the data that it produces. Unlike the eyebot,  it is powered as long as the cells in the car remain energized, so it will always log and monitor the health of the batteries.

\subsection{E\-xpert protocol}

The e-xpert device has a set protocol that it uses to communicate with other devices. It uses a RS232 connection over a 9 pin plug. This is a common way of communicating with external modules, and the eyebot has a serial port available to communicate with the e-xpert pro module. The module communicates using asynchronous communication. It automatically sends out updates at a rate of 1hz\cite{e_xpert}. These updates contain all the information that is recorded by the e-expert pro module. As the communication is asynchronous, this will happen automatically, even while the eyebot is not connected. This is not a problem as the e-xpert pro does not expect a response. This mode of operation is referred to as broadcast mode in the e-xpert documentation \cite{e_xpert}.

\begin{table}
\begin{center}
    \begin{tabular}{|l|l|}
        \hline
        Property & Value \\ \hline
        Baud Rate    & 2400 \\
        Data Bits    & 8    \\ 
        Stop Bits    & 1    \\ 
        Parity       & Even \\ 
        Flow Control & None \\
        \hline
    \end{tabular}
	\caption{Connection settings for Battery Monitor}
\end{center}
\end{table}


\subsubsection{Destination and Start Byte}
\label{sec:expertprotocol}

The message data that the module outputs is transmitted via serial. The first byte in the message is the start byte, in order to identify this start byte as the start byte it must be unique and never occur anywhere inside the payload. This is done by reserving the most significant bit (MSB) to be one only if it is a start or ending byte. The documentation refers to this bit as the IDHT (Identify Header Trailer ) bit. This does mean that there are only 7 bits available in each byte for transmitting data, but guarantees that the start and end of messages can be synchronized. As the first bit is a one due to the start byte being the header, the value of this byte is greater than 0x80. The rest of the bits in this first byte are the destination address. While communicating with a PC in broadcast mode, these bits can be ignored \cite{e_xpert}. The module also will not know where it is sending the byte, it is in broadcast mode, so there is no destination address. Thus the destination address bits will be 0, so the first header byte is always received as 0x80.

\subsubsection{Source}

The next byte transmitted is the source address. This byte is not a IDHT byte, so the MSB will always be 0. The device installed in the car, "e-xpert pro" will always set the source address as being 0. Combining this with the IDHT bit results in the second byte always being hex 0.

\subsubsection{Device ID}

The third byte in the message is the device ID. This is a unique number that identifies the type of equipment being used. This number is set by the manufacturer to distinguish different devices it it's product range. For the case of this product, the "e-xpert pro" the device id is 0x22.

\subsubsection{Message Identifier}

The e-xpert pro module transmits a variety of messages, which can be categorized into three groups. These different groups are handshake, commands and data. In broadcast mode, handshake and command messages are not required in order to extract information from the battery monitor. Table \ref{tab:tbsmessages} shows the hexadecimal values for different messages.

\begin{table}
\begin{center}
    \begin{tabular}{|l|l|}
        \hline
        Property & Value \\ \hline
        Battery Voltage   & 0x60 \\
        Battery Current    & 0x61    \\ 
        Amphours    & 0x62    \\ 
        Charge       &  0x64 \\ 
        Time Remaining & 0x65 \\
        \hline
    \end{tabular}
	\caption{Hexadecimal values of different TBS messages}
	\label{tab:tbsmessages}
\end{center}
\end{table}

\subsubsection{Data}

Following the message identifier is the actual payload of the message. This can take on various forms, but in the most simplistic sense is a number spread across a few bytes. Figure \ref{tbsdatafigure} shows the data layout for the battery voltage message. For more information about the different values, see "e-xpert pro communications interface specifications" \cite{e_xpert}.  Due to the MSB of each byte being reserved, it is only possible to transmit 7 bits of data inside a byte, any value that requires more than 7 bits to be represented must be transmitted across multiple bytes.  

\figuremacro{tbsdatafigure}{Structure of data payload}{Note that each byte transmitted only has seven usable bits of data, the 8th bit becomes the lowest bit in the next byte \cite{e_xpert}}

\subsubsection{Trailing Byte}

The last byte in the message is the "end of transfer" byte \cite{e_xpert}. The purpose of this byte is to signal that the message has been sent. Like the starting byte, the MSB of this byte is set to one. The rest of the bits in this byte are also set to one, to signify that this is the end byte, rather than the start byte. The value transmitted is 0xFF, and this is the only location in which 0xFF can appear.


\subsection {Design}

Like the GPS module, the battery monitor module was developed using C. This keeps in line with the design principles of making each component as simple as possible. Figure \ref{tbsflowchart} shows the program flow of this daemon.

\figuremacro{tbsflowchart}{Flow chart of the battery monitor daemon}{}

\subsubsection{Network protocol}

Table \ref{tab:tbsprotocol} shows the protocol for the battery monitor message when transmitted over the network. As the design uses a binary format file like the GPS module, this format is impossible to read natively, but more efficient on space. There is no need for these values to be read on the network layer anyway. All values are transmitted in big endian network order.

\begin{table}
    \begin{tabular}{*{7}{|l}|}
        \hline
        54 &  42 & 53 & 0 & 0  & 4e1fe968  &         00000000   \\ \hline \hline
        \multicolumn{5}{|c|}{ Header} & \multicolumn {2}{|c|}{Data Payload} \\ \hline
	\multicolumn{3}{|c|}{Filter} & Minor Filter & Revision Number &Time(s) & Time remaining (s) \\
        \hline
    \end{tabular}
%\end{table}
%
%\begin{table}
    \begin{tabular}{*{6}{|l}|}
        \hline
 42b40000 & 43710000 & c150000 &  00000000 & 00000000 & 00000000 \\ \hline \hline
         \multicolumn {6}{|c|}{Data Payload} \\ \hline
	  Charge(\%) & Voltage(V) & Current(A) & AmpHours(Ah) & Temperature(C) & Status \\
        \hline
    \end{tabular}
	\caption{Network protocol for Battery Monitor daemon}
	\label{tab:tbsprotocol}
\end{table}



% ---------------------------------------------------------------------------
%: ----------------------- end of thesis sub-document ------------------------
% ---------------------------------------------------------------------------



%: ----------------------- contents from here ------------------------

\section{Trip Meter}



% ---------------------------------------------------------------------------
%: ----------------------- end of thesis sub-document ------------------------
% ---------------------------------------------------------------------------



%: ----------------------- contents from here ------------------------

\section{Interial Measurement Unit Display Panel}

Developed in section 


\figuremacro{imu}{IMU display panel}{This figure shows the information output by the IMU daemon }



% ---------------------------------------------------------------------------
%: ----------------------- end of thesis sub-document ------------------------
% ---------------------------------------------------------------------------





% ---------------------------------------------------------------------------
%: ----------------------- end of thesis sub-document ------------------------
% ---------------------------------------------------------------------------

	
% this file is called up by thesis.tex
% content in this file will be fed into the main document

%: ----------------------- name of chapter  -------------------------
\chapter{Interface} % top level followed by section, subsection


%: ----------------------- paths to graphics ------------------------

% change according to folder and file names
\ifpdf
    \graphicspath{{6/figures/PNG/}{6/figures/PDF/}{6/figures/}}
\else
    \graphicspath{{6/figures/EPS/}{6/figures/}}
\fi

%: ----------------------- contents from here ------------------------


\section{Layout}

An important aspect of user interfaces is that they must be visually impressive [GO CITE?]. This keeps the user interested in the product, and helps new users want to learn how to operate the device. The work done in the windowing toolkit allowed a much more visually impressive layout to be developed. This was achieved through the use of transparencies in order to overlay different elements on top of each other.

\subsection{Background}

As the device does not have enough power to dynamically render any sort of complex images or motions, a pre-rendered static background was used. This background was designed on another pc, and converted for use on the embedded system. It consists of two main regions. The first area is the main display region. This section is the largest of the space and exists so the current data being displayed can be laid over it. The second region is the navigation bar region. This area is at the bottom of the image, it uses a more uniform texture in order to contrast the main data display area above it. The background itself is rendered in black and white. The black and white styling allows important information to contrast against it easily. This helps highlight the information and user interface actions present to the user, and avoids a simple single-colour background on the user interface.

\subsection{Navigation Model}

Figure \ref{uinavigation} shows the basic routes throughout the user interface. This model is a generic tree, that can have as many leafs as possible. When the user wishes to navigate to a new panel, a pointer to that new panel is required. Thus as-long as the panel knows of the panel it wishes to transition to, any panel can be made active from any other panel. This allows the nested elements, such as syslog and status to navigate to any of the higher level panels easily as-well. These routes are not shown on the diagram, as there would be too many pathways to follow. If a panel returns a NULL pointer, the system will navigate to the App Panel to avoid a crash. This also provides handy functionality for traversing back to the top-most panel in the tree, as any leaf panel can request to transition to NULL and the App Panel will be displayed.


\figuremacro{uinavigation}{Basic Navigation of the interface}{The basic navigation routes through the user interface, does not show "quickbar" navigation routes}



%: ----------------------- contents from here ------------------------

\section{Overview Panel}

To facilitate easy navigation a panel was designed that shows all the different panels available. This panel is shown in figure \ref{app}. It contains 8 different aspects of the program, with room available for 12. This will allow the system to be extended in the future. 

\figuremacro{app}{Panel showing other panels}{Shortcuts to different aspects of the program. (Battery, Maps, Trip Meter, Accelerometer, Arduino, Savings, About, Options)}



% ---------------------------------------------------------------------------
%: ----------------------- end of thesis sub-document ------------------------
% ---------------------------------------------------------------------------



%: ----------------------- contents from here ------------------------

\section{Battery}

The battery panel is the main display panel used in the user interface. It is the first panel displayed to the user when the system turns on.  It displays five important pieces of information to the user operating the vehicle. This panel listens to both the battery (TBS) and gps messages. It requires the battery messages to display the battery voltage, current and charge to the user. These are displayed using the digitelements mentioned in \ref{sec:digitelement}. It also uses a custom charge element class that exists outside of the framework to draw a battery on the center of the screen. This green bar on this battery will decrease in proportion to the charge remaining in the car, providing a quick visual indicator for how much charge is still stored. 

This panel also displays the speed in km/h, which is why it has to register to the gps module. The speed is included so the user does not have to interact with the screen while driving. This panels main purpose was to provide a quick overview to data points that are immediately of concern to the driver. 

Displayed in the top left corner is a rudimentary calculation of how much distance is remaining in the car. Range tests conducted by the REV team indicated that this distance was 80km from a full charge. To provide some security while driving, the interface assumes that the max distance the car will travel on a full charge is 70km. It uses this to calculate the distance remaining according to the formula \ref{eq:distanceremaining}

\begin{equation}
\label{eq:distanceremaining}
70=10
\end{equation}

\figuremacro{batt}{Battery state panel}{Screenshot of the battery state panel, showing the voltage, current, charge, speed (via gps) and distance remaining}



% ---------------------------------------------------------------------------
%: ----------------------- end of thesis sub-document ------------------------
% ---------------------------------------------------------------------------



%: ----------------------- contents from here ------------------------

\section{Maps}

A common, yet useful driver aid is displaying the map of the current location. This is what the maps panel does. It listens to the GPS messages to determine the position of the car. A screen-shot of this panel is shown in Figure \ref{map}. This panel uses the full area to display the current location of the vehicle on the screen. The actual position of the vehicle is centered on the middle of the screen, allowing the driver to view streets and landmarks relative to his current position. The panel features several levels of zoom, controlled by the bar on the top right corner on the screen. The plus button will zoom in, and move the slider to the right, while the minus button will zoom out, and move the slider to the left. The system supports various levels of zoom, being able to display a few buildings relative to the car or being able to display the surrounding suburbs.

\figuremacro{map}{Map display panel}{Screen-shot of the map display panel, showing the map of the current location and the map controls visible}

\subsection{Map Data}

In order to display the map images, the map data must first be obtained and stored. There are many possible methods for doing this, ranging from creating the maps as needed, or downloading them from external services and storing them in a cache. The relatively low CPU power of the eye-bot m6 makes rendering the maps on-the-fly a undesirable prospect. Open source map data for the entire planet results in a file that is 18GB in size \cite{planet_osm}. This file is much too large to store locally, and this file is actually storing a compressed version of the data. Even if it were possible for the eye-bot to store this file, via the use of external storage, it would be a large strain on the CPU to convert the street level data into viewable maps. It would also require many other pieces of software to be installed on the device, making it much more complicated to manage.

Another possibility is to use an existing Internet based map server and download the map imagery as needed. The has several downsides. Foremost it requires an Internet connection whenever to display the maps. The system does have a 3G connection installed, but this cannot be considered a dependable communication channel. It is highly likely to drop out, and is limited in coverage to the areas in which it has reception. Even without these issues, most map-servers do not allow you to download maps in bulk, as this violates their usage policies \cite{tile_usage_policy}. This makes this method undesirable as it is not suit-able for downloading maps in bulk, or as needed. Attempting to pre-download all these maps using a non-3g link would result in a violation of the usage policy.

The method chosen to obtain the map data was to pre-create create the map data using a more powerful machine. A map-server was setup and loaded with all the street data for the oceania region. For more information on the map-server setup see appendix \ref{app:mapserver}. This method overcomes the problems of the previously mentioned methods. All the processing is done on the much faster machine in advance. The area processed in advance is defined by the properties in table \ref{tab:mapbbox}. This area is depicted by the image shown in Figure \ref{mapperth}. The expanse of these maps covers all of metropolitan Perth. In future more maps can easily be processed, however this will result in more storage space being required. The current settings are a good balance between storage and expanse of data. This is because the map data will be less useful outside of the city, and the car is typically not driven any further than the pre-rendered maps. The current space usage of the map data is given in table \ref{tab:mapdata}. As this is much much larger than the internal 16MB of flash storage the eyebot has, it must be stored on external storage. Luckily, storage is now cheap, and 8GB usb thumb drive has enough performance and space to store and serve the image data.

\begin{table}[htdp]
\begin{center}
	
   \begin{tabular}{|l|l|}
        \hline
        Property           & Value           \\ \hline
        Minimum Zoom Level & 11              \\ 
        Maximum Zoom Level & 18              \\ 
        Top Left           & 115.687,-31.71  \\ 
        Bottom Right       & 116.508,-32.253 \\
        \hline
    \end{tabular}
	\caption[Properties of the pre-rendered map data]{The properties of the pre-rendered map data}
	\label{tab:mapbbox} 
\end{center}
\end{table}

\figuremacro{mapperth}{Pre-rendered map size}{This figure shows the area that has been pre-rendered for use in the map panel display}

\begin{table}[htdp]
\begin{center}
    \begin{tabular}{|l|l|}
        \hline
        Property              & Value                             \\ \hline
        Zoom Levels           & 7                                 \\ 
        Number of Files       & 435,456 \\ 
        Total Size            & 580.2 MB                          \\
        File Resolution            & 256*256 pixels                          \\  
        File Format            & Palette PNG                         \\ 
        Time taken to process & Approximately 1 Day                          \\
        \hline
    \end{tabular}
	\caption[File statistics of map data]{File statistics of map data}
	\label{tab:mapdata}
\end{center}
\end{table}

\subsection{Tiling}

As mentioned previously, the map data is much larger than the internal storage of the eye-bot. It is also much much larger than the operating systems 64MB of ram. This makes it impossible for the entire map data to be loaded inside any program to be displayed to the user. In order to overcome this limitation, the map panel only loads a small subset of the map data at a time, and displays it using a method called tiling. Rather than have one giant map that displays all of the information, and sliding this around to view the correct part, this method divides the map into many smaller maps. These smaller maps are referred to as tiles. Each tile consists of a 256 pixels wide by 256 pixels high image. These images are combined in order to display the current location, as seen in figure \ref{tiles}. By loading the adjacent tiles in all directions the screen will always have data to display.

\figuremacro{tiles}{Tiling Maps}{This figure shows the tiles that are loaded, labeled 0-8, and the area of the screen that is able to view them.}

\subsubsection{Converting GPS Co-ordinates}

The GPS outputs the current position in latitude/longitude format. This format must be converted so that the correct tiles may be loaded. To do this the GPS co-ordinates are converted to grid based co-ordinates using a mercator projection \cite{slippy_map_tilenames}. The equations to convert to this format are given by \ref{eq:gpstotiles}. It is possible to convert from a grid co-ordinate back to a GPS co-ordinate using the equations given by \ref{eq:tiletogps}. This of course will only be able to return the top left corner of the tile in GPS co-ordinates, and not the original position.

\begin{equation}
\label{eq:gpstotiles}
70=10
\end{equation}
\begin{equation}
\label{eq:tilestogps}
70=10
\end{equation}

With the GPS co-ordinate converted it is possible to locate the tile to display to the user. Rather than utilize lookup files to locate the tile, a specific folder structure is used to instantaneously load the file. This is done by storing the files in the format \emph{/zoom/tilex/tiley.png}. This allows the system to directly open the file and display it. Thus even if the entire world was loaded in the maps folder, for many different zoom levels, the display time of the current location would not change.

% ---------------------------------------------------------------------------
%: ----------------------- end of thesis sub-document ------------------------
% ---------------------------------------------------------------------------



%: ----------------------- contents from here ------------------------

\section{Trip Meter}

A useful driver aid that is common on vehicles is that of a trip meter. Traditionally this component records the distance the car has traveled since the trip meter was set. This functionality is usually a result of the simple systems in place, and is tied to the revolutions of the wheels on the vehicle. As this system has more hardware at it's disposal, the trip meter can implement more functionality than a standard trip meter, making it much more useful in examining the performance of the car. Figure \ref{trip} shows the trip meter panel being displayed on the screen.

\figuremacro{trip}{Trip Meter Panel}{This figure shows two independent trip meters and the best record speed data}

The first unique point of this trip meter, is that it contains two independent meters. This is useful as it allows the driver to evaluate the statistics of two overlapping trips. One trip meter can be used to record the distance traveled since the car was last charged, while the other can be used to record the distance traveled in the last week or month. This independence allows the operator to decide how best to use the trip meter data, resulting in a high level of flexibility. In figure \ref{trip} the trip meters are located on the left, sitting above each other.

Each trip meter records the distance traveled, the time elapsed since the meter was started, the time the car has been moving since the meter was started and calculations based on the elapsed and moving time. These statistics are displayed live to the user, but are not logged, as the logging functionality  is taken car of by a different component in the system. The trip meter panel also displays the current moving speed in the top right. Below the moving speed are the best run records in seconds. This allows the driver to have quick feedback as to how the car is performing, without having to do lots of processing on logged data.

\subsection{Distance Driven}
\label{sec:distancedriven}
The most important part of a trip meter, is the distance that the meter has recorded. This is shown on figure \ref{trip} at the bottom of each trip meter. The meter stores the distance driven internally as a double length floating point number, but displays it on the screen as a rounded integer. This is done in order improve precision for later calculations as other values will depend on the distance that has been driven. Equation \ref{eq:distancedriven} shows the the distance is calculated based on the GPS position.

\begin{align}
\label{eq:distancedriven}
\mathrm{Distance} &=\mathrm{Distance}_\mathrm{lastrun} + \mathrm{Speed}(\mathrm{Time}_\mathrm{now}-\mathrm{Time}_\mathrm{lastrun})
\end{align}

The method for working out is a continuous function that is based on the last known distance the car has traveled. This method was chosen as it does not require any information other than the last time the formula was run, and the last distance calculated. This also makes the trip meter flexible in that the distance calculation does not need to be processed at exact intervals. If messages are dropped for whatever reason, the calculation will still take place, though it will not be as accurate as it could be. The calculation will be able to cope with fluctuations in the message timing, and can adapt to the speed of the GPS being used.

The downside of this method, is that big changes in time can cause problems with the calculation. If the signal drops out for an extended period of time, such as going through a tunnel, the calculation in \ref{eq:distancedriven} would have a big margin for error. In order to prevent this, the trip meter will ignore large time differences. If two calculations are over 10 seconds apart, the result will not be trusted, and not be used in the calculation. This prevents GPS signal loss from having an adverse effect on the trip meter calculations, but does impose a limit on the trip meter.

The limitation of this method of calculating the distance driven is that it relies on the GPS messages being sent to it. If the GPS signal is lost, the distance driven will not be increased. This is a limitation imposed by the use of the GPS sensor, and cannot be avoided, as attempting to guess the distance driven while the GPS signal has been lost has a very high probability of being incorrect.

\subsection{Time Elapsed}

Another variable displayed on the trip panel is the time elapsed. This is simply the time elapsed since the trip meter was last reset. This time increments even while the GPS signal has been lost, performing a stop-watch like action on the trip. While it may seem natural to just record the time that the trip meter was started and subtract it from the current system time, this would lead to problems during the system startup. The time needs to increment even while the GPS is connecting, and must be resistant to changes in the systems internal clock. Thus the time is calculated similar to section \ref{sec:distancedriven}. The formula used to calculate the elapsed time is given in equation \ref{eq:timeelapsed}. The time elapsed is displayed as the highest element of the trip meter in Figure \ref{trip}

\begin{align}
\label{eq:timeelapsed}
\mathrm{TimeElapsed} &=\mathrm{TimeElapsed}_\mathrm{lastrun} + (\mathrm{Time}_\mathrm{now}-\mathrm{Time}_\mathrm{lastrun})
\end{align}

Much like the distance, this method of calculation depends on the last known values. This means it does not matter at the actual time the system trip started once the timer has been running. This makes it resistant to changes in the system time, and thus makes the timer more robust. 
This timer records the time elapsed on the nanosecond level, as it is used elsewhere in calculations. For display, the timer converts these values into the traditional hours, minutes, seconds format that is easy for the operator to read.

\subsection{ Moving Time}

An aspect of the cars telemetry that would be interesting to the driver is the cars moving time. This is defined as the time in which the car has spent in motion. The main use of this data point is to contrast it against the elapsed time, to highlight how long the car has spent sitting still in traffic. This variable is also useful to record for future calculations, such as working out the average speed of the trip. This element is calculated according to equation \ref{eq:timeelapsed}, except that it will not update if the current speed of the car is 0~km/h. As such this element requires the speed of the car to be processed, so it cannot be calculated when the GPS signal is lost. This fits in with the functionality defined in section \ref{sec:distancedriven}, as the trip meter will not update the moving time or distance driven if the GPS signal is lost. The moving time is displayed below the elasped time and above the distance driven in Figure \ref{trip}

\subsection{Average Speed}

When reviewing the trip meter data, it is useful to know the average speed the car was traveling during the trip. Having this information allows the driver to better understand the characteristics of the drive. This element is also easy to calculate, as the time elapsed and the distance driven are already available. Equation \ref{eq:averagespeed} shows the formula used to calculate the average speed. The calculated value is displayed to the right of the elapsed time in figure \ref{trip}.

\begin{align}
\label{eq:averagespeed}
\mathrm{Average Speed} &=\frac{\mathrm{Distance Driven}}{\mathrm{Time Elapsed}}
\end{align}

This value is calculated whenever the distance driven or elapsed time values are updated. As this value is using two calculated values, it does not need to worry about discrepancies in time or the loss of the GPS signal.

\subsection{Average Moving Speed}

The average moving speed is like the average speed. The only difference is that it uses the moving time to calculate the speed, rather than the elapsed time. This is done using the same equation \ref{eq:averagespeed}, only substituting "Time Elapsed" for "Moving Time". This value provides the average speed of the car when it was actually being driven, thus ignoring time spent waiting in traffic.

\subsection{Reset}

The final functionality of each independent trip meter is the reset button. This button resets the trip meter it is attached to. The distance driven, moving and elapsed time counters will all display zero, and the average speed calculators will display zero. As each trip meter is independent, one can be reset without affecting the other. To reset the trip meters, the driver just has to press the reset button, located below the average moving time and to the right of the distance driven in Figure \ref{trip}.

\subsection{Current Speed}

The trip meters provide statistics on where the car has been driven, but do not provide much insight into the instantaneous speed of the vehicle. As this variable is already being used in calculations it is trivial to display it to the driver. This is displayed using a simple digit element, and appears in the top right corner of Figure \ref{trip}.


% ---------------------------------------------------------------------------
%: ----------------------- end of thesis sub-document ------------------------
% ---------------------------------------------------------------------------



%: ----------------------- contents from here ------------------------

\section{Interial Measurement Unit Display Panel}

Developed in section 


\figuremacro{imu}{IMU display panel}{This figure shows the information output by the IMU daemon }



% ---------------------------------------------------------------------------
%: ----------------------- end of thesis sub-document ------------------------
% ---------------------------------------------------------------------------



%: ----------------------- contents from here ------------------------

\section{Trip Meter}



% ---------------------------------------------------------------------------
%: ----------------------- end of thesis sub-document ------------------------
% ---------------------------------------------------------------------------



%: ----------------------- contents from here ------------------------

\section{Trip Meter}



% ---------------------------------------------------------------------------
%: ----------------------- end of thesis sub-document ------------------------
% ---------------------------------------------------------------------------



%: ----------------------- contents from here ------------------------

\section{About}

% ---------------------------------------------------------------------------
%: ----------------------- end of thesis sub-document ------------------------
% ---------------------------------------------------------------------------




%: ----------------------- contents from here ------------------------

\section{Settings}

The previous features have all discussed elements of the interface that are useful to any general driver. These are features that are useful while operating the vehicle. There also exists another user that has a different set of requirements to fulfill. This user is the maintainer of the car. This user will need access to various aspects for debugging and development purposes, including the ability to easily examine and transfer the internal logs of the system. In order to help simplify the navigation of the program, these elements are located in the settings screen. A screen-shot of the settings screen is shown in figure \ref{options}.

\figuremacro{options}{Options Panel}{This figure shows the options panel, which displays various utilities that are needed by the maintainer of the vehicle.}


% ---------------------------------------------------------------------------
%: ----------------------- end of thesis sub-document ------------------------
% ---------------------------------------------------------------------------



%: ----------------------- contents from here ------------------------

\section{Debug}

The top left most button in figure \ref{options} will navigate to the debug panel. This panel is shown in figure \ref{messages}. The purpose of this panel is to display the messages that are being received from the daemon programs, see section \ref{sec:design}. This panel displays the byte level contents of each message to the user as it is received. This is used to help facilitate debugging, as the writer of the daemon is able to see exactly what is being received in the interface.

\figuremacro{messages}{Debug Messages Panel}{This figure shows the Debug panel, it allows the user to filter which messages are displayed and alternates the colour between each received message}

This display panel has various buttons that exist on the top of the screen. These buttons allow known messages to be filtered out, enabling specific debugging to be undertaken. As discussed earlier, the first 3 bytes of a message are used to filter different kinds of messages. When a button is enabled, it will allow those messages to be drawn on the screen. Multiple combinations of messages can be combined to display any combination. The buttons labeled, TBS,GPS,IMU and DIO will filter the ASCII values of the first three bytes of each message. The all button will ignore any filtering, and will allow any message to be displayed, including ones that may not have been developed or seen before.

This panel uses a ConsoleElement to display it's data as seen in section \ref{sec:consoleelement}. This allows new messages to be displayed on the bottom of the screen, and travel upwards as they become outdated. It also provides a clear distinction between messages, as seen in figure \ref{messages}. The messages alternate between red and white as they are displayed, making it easy to identify the trailing portions of each message.

\section{Network Status Display}

In order to verify that the system is working correctly, it is helpful to view the messages being sent as in the Debug Panel. This panel can provide too much information, and requires the user to understand hex in order to verify that each message is being received. In order to simplify this, a panel was developed that displays the status of each different type of message. This uses the simple boolean display element previously mentioned in section \ref{sec:diopanel}. This element will display green whenever messages are currently being received by the system. If no message is received in the previous 10 seconds, the marker will transition to red in order to illustrate this. This allows a quick view as to the health of the system, and highlights components that need to have further inspection.

\figuremacro{netstat}{Daemon Status Panel}{This figure shows the status of the various daemons, green indicates the daemon is active, red indicates a timeout}

\section{System Logs Display}

As mentioned in section \ref{sec:dmesglogging}, all aspects of this system report errors by logging them to the syslog application. This prints out the warnings and errors created to \emph{/var/log/messages}. This panel loads and displays any information logged into this file. This information may come from any of the daemons mentioned in section \ref{sec:design}, the interface itself, or even other systems that are running on the Linux installation.

\figuremacro{dmesg}{Syslog Panel}{This figure shows the System Log panel, this displays the most recent entries in \emph{/var/log/messages}}

An important aspect of this panel is that it only displays the most recent information that has been logged. This makes it compatible with the ConsoleElement described in section \ref{sec:consoleelement}. The Console Element will automatically take care of displaying the last pieces of information it has received. This would produce the desired functionality, however further optimization needed to be undertaken.

The file \emph{/var/log/messages}, is logged to by all applications since the system started. Thus it grows in size the longer the system has been operational. Without optimizing the way this file is read, it would take O(n) time to display the newest entries in the file on the Console Element, where n is the number of entries in the file. This is not acceptable as the time taken to display the logs will increase the longer the application is running. In order to fix this problem, the file is read from the end of the file minus a 2048 byte offset.  The maximum width that can be displayed is 60 characters, and the maximum possible lines that can be shown on the screen is 20. This yields the result of 1200 possible characters as seen in equation \ref{eq:maximumdisplaycharacters}. This size is roughly doubled to 2048, in order to ensure more than enough lines are read for each update. This method of reading the last 2048 characters and finding the unique lines inside them is not affected by the file size, so the time taken to update the display is now O(1). The update is not dependent on the number of lines in \emph{/var/log/messages}

\begin{align}
\label{eq:maximumdisplaycharacters}
\mathrm{Maximum\ Characters} &= \mathrm{Characters\ per\ line}* \mathrm{Number\ of \ Lines} \\
&= 20*60 \notag \\
&=1200 \notag 
\end{align}

As the actual log file may be updated very quickly the screen must not update automatically. This is to stop the screen transitioning while the user may be reading something on it. As such the panel does not support any automatically updating functionality. In order to update the screen, a system like that in section \ref{sec:screenabout} is used. A invisible button is created to span the whole display. When pressed, this button will read the \emph{/var/log/messages} file and update the screen accordingly. This allows for the maximum utilization of space possible, as no space needs to be created to display a refresh button. The messages will also be updated when this panel becomes the active panel displayed on the screen.


% ---------------------------------------------------------------------------
%: ----------------------- end of thesis sub-document ------------------------
% ---------------------------------------------------------------------------



%: ----------------------- contents from here ------------------------

\section{Copy Log Files}

Discussed in section \ref{sec:filelogger} is a utility that records the telemetry data locally. In order to make the files created by this logger easier to access, the user interface provides a mechanism for them to be copied onto an external usb drive. Figure \ref{copy} shows the display asking for user confirmation to copy the log files. If the user confirms the action, the program will navigate to the directory in which the log files are stored, and copy them byte by byte to the external drive. Once completed the screen will display the figure \ref{copycompleted}. 
\figuremacro{copy}{Copy Popup}{Popup asking the user if they wish to copy log files to the an external drive}

If an error occurs, the system will log the results to the in-built system logger, syslog. This will allow the results to be viewable at \emph{/var/log/messages}, and thus can be viewed by the Syslog panel. While copying, the interface will display the progress of the action. It does this by indicating the amount of files to be copied, and the current file it is copying across. This is implemented so the user can be sure the system has not hung or encountered another error.

\figuremacro{copycompleted}{Copy Completed Popup}{This popup confirms that all the log files have been correctly copied to the usb drive}

\section{Delete Log Files}

In order to prevent the storage media from filling up, and to help organize the files when they are copied over the interface provides a mechanism to delete all the log files. Figure \ref{delete} shows a popup asking the user if they wish to delete the files. This functionality functions much the same way as copy, except that it does not require an external drive. It will simply open the folder and delete the files one by one. This action will also display its progress, and will provide the user with a message once the action is complete.

\figuremacro{delete}{Delete Log Files Popup}{This asks the user if they want to delete all the log files stored on the system}

% ---------------------------------------------------------------------------
%: ----------------------- end of thesis sub-document ------------------------
% ---------------------------------------------------------------------------


\section{Exit Program}

In order to allow other programs that use the various hardware, such as the screen or the touch driver, to run the current application must be exited. The system facilitates this be placing an exit action on the options panel. When this button is pressed, the popup shown in figure \ref{exit} will be displayed. This popup confirms that the user wishes to exit the application. If the user agrees, the program will exit and the screen will be set to black. If the user cancels, the system will return to the previously displayed panel, allowing normal operation to resume.

\figuremacro{exit}{Exit Popup}{This popup confirms that the user wants to exit the running application}




% ---------------------------------------------------------------------------
%: ----------------------- end of thesis sub-document ------------------------
% ---------------------------------------------------------------------------


% this file is called up by thesis.tex
% content in this file will be fed into the main document

%: ----------------------- name of chapter  -------------------------
\chapter{Interface} % top level followed by section, subsection


%: ----------------------- paths to graphics ------------------------

% change according to folder and file names
\ifpdf
    \graphicspath{{6/figures/PNG/}{6/figures/PDF/}{6/figures/}}
\else
    \graphicspath{{6/figures/EPS/}{6/figures/}}
\fi

%: ----------------------- contents from here ------------------------


\section{Layout}

An important aspect of user interfaces is that they must be visually impressive [GO CITE?]. This keeps the user interested in the product, and helps new users want to learn how to operate the device. The work done in the windowing toolkit allowed a much more visually impressive layout to be developed. This was achieved through the use of transparencies in order to overlay different elements on top of each other.

\subsection{Background}

As the device does not have enough power to dynamically render any sort of complex images or motions, a pre-rendered static background was used. This background was designed on another pc, and converted for use on the embedded system. It consists of two main regions. The first area is the main display region. This section is the largest of the space and exists so the current data being displayed can be laid over it. The second region is the navigation bar region. This area is at the bottom of the image, it uses a more uniform texture in order to contrast the main data display area above it. The background itself is rendered in black and white. The black and white styling allows important information to contrast against it easily. This helps highlight the information and user interface actions present to the user, and avoids a simple single-colour background on the user interface.

\subsection{Navigation Model}

Figure \ref{uinavigation} shows the basic routes throughout the user interface. This model is a generic tree, that can have as many leafs as possible. When the user wishes to navigate to a new panel, a pointer to that new panel is required. Thus as-long as the panel knows of the panel it wishes to transition to, any panel can be made active from any other panel. This allows the nested elements, such as syslog and status to navigate to any of the higher level panels easily as-well. These routes are not shown on the diagram, as there would be too many pathways to follow. If a panel returns a NULL pointer, the system will navigate to the App Panel to avoid a crash. This also provides handy functionality for traversing back to the top-most panel in the tree, as any leaf panel can request to transition to NULL and the App Panel will be displayed.


\figuremacro{uinavigation}{Basic Navigation of the interface}{The basic navigation routes through the user interface, does not show "quickbar" navigation routes}



%: ----------------------- contents from here ------------------------

\section{Overview Panel}

To facilitate easy navigation a panel was designed that shows all the different panels available. This panel is shown in figure \ref{app}. It contains 8 different aspects of the program, with room available for 12. This will allow the system to be extended in the future. 

\figuremacro{app}{Panel showing other panels}{Shortcuts to different aspects of the program. (Battery, Maps, Trip Meter, Accelerometer, Arduino, Savings, About, Options)}



% ---------------------------------------------------------------------------
%: ----------------------- end of thesis sub-document ------------------------
% ---------------------------------------------------------------------------



%: ----------------------- contents from here ------------------------

\section{Battery}

The battery panel is the main display panel used in the user interface. It is the first panel displayed to the user when the system turns on.  It displays five important pieces of information to the user operating the vehicle. This panel listens to both the battery (TBS) and gps messages. It requires the battery messages to display the battery voltage, current and charge to the user. These are displayed using the digitelements mentioned in \ref{sec:digitelement}. It also uses a custom charge element class that exists outside of the framework to draw a battery on the center of the screen. This green bar on this battery will decrease in proportion to the charge remaining in the car, providing a quick visual indicator for how much charge is still stored. 

This panel also displays the speed in km/h, which is why it has to register to the gps module. The speed is included so the user does not have to interact with the screen while driving. This panels main purpose was to provide a quick overview to data points that are immediately of concern to the driver. 

Displayed in the top left corner is a rudimentary calculation of how much distance is remaining in the car. Range tests conducted by the REV team indicated that this distance was 80km from a full charge. To provide some security while driving, the interface assumes that the max distance the car will travel on a full charge is 70km. It uses this to calculate the distance remaining according to the formula \ref{eq:distanceremaining}

\begin{equation}
\label{eq:distanceremaining}
70=10
\end{equation}

\figuremacro{batt}{Battery state panel}{Screenshot of the battery state panel, showing the voltage, current, charge, speed (via gps) and distance remaining}



% ---------------------------------------------------------------------------
%: ----------------------- end of thesis sub-document ------------------------
% ---------------------------------------------------------------------------



%: ----------------------- contents from here ------------------------

\section{Maps}

A common, yet useful driver aid is displaying the map of the current location. This is what the maps panel does. It listens to the GPS messages to determine the position of the car. A screen-shot of this panel is shown in Figure \ref{map}. This panel uses the full area to display the current location of the vehicle on the screen. The actual position of the vehicle is centered on the middle of the screen, allowing the driver to view streets and landmarks relative to his current position. The panel features several levels of zoom, controlled by the bar on the top right corner on the screen. The plus button will zoom in, and move the slider to the right, while the minus button will zoom out, and move the slider to the left. The system supports various levels of zoom, being able to display a few buildings relative to the car or being able to display the surrounding suburbs.

\figuremacro{map}{Map display panel}{Screen-shot of the map display panel, showing the map of the current location and the map controls visible}

\subsection{Map Data}

In order to display the map images, the map data must first be obtained and stored. There are many possible methods for doing this, ranging from creating the maps as needed, or downloading them from external services and storing them in a cache. The relatively low CPU power of the eye-bot m6 makes rendering the maps on-the-fly a undesirable prospect. Open source map data for the entire planet results in a file that is 18GB in size \cite{planet_osm}. This file is much too large to store locally, and this file is actually storing a compressed version of the data. Even if it were possible for the eye-bot to store this file, via the use of external storage, it would be a large strain on the CPU to convert the street level data into viewable maps. It would also require many other pieces of software to be installed on the device, making it much more complicated to manage.

Another possibility is to use an existing Internet based map server and download the map imagery as needed. The has several downsides. Foremost it requires an Internet connection whenever to display the maps. The system does have a 3G connection installed, but this cannot be considered a dependable communication channel. It is highly likely to drop out, and is limited in coverage to the areas in which it has reception. Even without these issues, most map-servers do not allow you to download maps in bulk, as this violates their usage policies \cite{tile_usage_policy}. This makes this method undesirable as it is not suit-able for downloading maps in bulk, or as needed. Attempting to pre-download all these maps using a non-3g link would result in a violation of the usage policy.

The method chosen to obtain the map data was to pre-create create the map data using a more powerful machine. A map-server was setup and loaded with all the street data for the oceania region. For more information on the map-server setup see appendix \ref{app:mapserver}. This method overcomes the problems of the previously mentioned methods. All the processing is done on the much faster machine in advance. The area processed in advance is defined by the properties in table \ref{tab:mapbbox}. This area is depicted by the image shown in Figure \ref{mapperth}. The expanse of these maps covers all of metropolitan Perth. In future more maps can easily be processed, however this will result in more storage space being required. The current settings are a good balance between storage and expanse of data. This is because the map data will be less useful outside of the city, and the car is typically not driven any further than the pre-rendered maps. The current space usage of the map data is given in table \ref{tab:mapdata}. As this is much much larger than the internal 16MB of flash storage the eyebot has, it must be stored on external storage. Luckily, storage is now cheap, and 8GB usb thumb drive has enough performance and space to store and serve the image data.

\begin{table}[htdp]
\begin{center}
	
   \begin{tabular}{|l|l|}
        \hline
        Property           & Value           \\ \hline
        Minimum Zoom Level & 11              \\ 
        Maximum Zoom Level & 18              \\ 
        Top Left           & 115.687,-31.71  \\ 
        Bottom Right       & 116.508,-32.253 \\
        \hline
    \end{tabular}
	\caption[Properties of the pre-rendered map data]{The properties of the pre-rendered map data}
	\label{tab:mapbbox} 
\end{center}
\end{table}

\figuremacro{mapperth}{Pre-rendered map size}{This figure shows the area that has been pre-rendered for use in the map panel display}

\begin{table}[htdp]
\begin{center}
    \begin{tabular}{|l|l|}
        \hline
        Property              & Value                             \\ \hline
        Zoom Levels           & 7                                 \\ 
        Number of Files       & 435,456 \\ 
        Total Size            & 580.2 MB                          \\
        File Resolution            & 256*256 pixels                          \\  
        File Format            & Palette PNG                         \\ 
        Time taken to process & Approximately 1 Day                          \\
        \hline
    \end{tabular}
	\caption[File statistics of map data]{File statistics of map data}
	\label{tab:mapdata}
\end{center}
\end{table}

\subsection{Tiling}

As mentioned previously, the map data is much larger than the internal storage of the eye-bot. It is also much much larger than the operating systems 64MB of ram. This makes it impossible for the entire map data to be loaded inside any program to be displayed to the user. In order to overcome this limitation, the map panel only loads a small subset of the map data at a time, and displays it using a method called tiling. Rather than have one giant map that displays all of the information, and sliding this around to view the correct part, this method divides the map into many smaller maps. These smaller maps are referred to as tiles. Each tile consists of a 256 pixels wide by 256 pixels high image. These images are combined in order to display the current location, as seen in figure \ref{tiles}. By loading the adjacent tiles in all directions the screen will always have data to display.

\figuremacro{tiles}{Tiling Maps}{This figure shows the tiles that are loaded, labeled 0-8, and the area of the screen that is able to view them.}

\subsubsection{Converting GPS Co-ordinates}

The GPS outputs the current position in latitude/longitude format. This format must be converted so that the correct tiles may be loaded. To do this the GPS co-ordinates are converted to grid based co-ordinates using a mercator projection \cite{slippy_map_tilenames}. The equations to convert to this format are given by \ref{eq:gpstotiles}. It is possible to convert from a grid co-ordinate back to a GPS co-ordinate using the equations given by \ref{eq:tiletogps}. This of course will only be able to return the top left corner of the tile in GPS co-ordinates, and not the original position.

\begin{equation}
\label{eq:gpstotiles}
70=10
\end{equation}
\begin{equation}
\label{eq:tilestogps}
70=10
\end{equation}

With the GPS co-ordinate converted it is possible to locate the tile to display to the user. Rather than utilize lookup files to locate the tile, a specific folder structure is used to instantaneously load the file. This is done by storing the files in the format \emph{/zoom/tilex/tiley.png}. This allows the system to directly open the file and display it. Thus even if the entire world was loaded in the maps folder, for many different zoom levels, the display time of the current location would not change.

% ---------------------------------------------------------------------------
%: ----------------------- end of thesis sub-document ------------------------
% ---------------------------------------------------------------------------



%: ----------------------- contents from here ------------------------

\section{Trip Meter}

A useful driver aid that is common on vehicles is that of a trip meter. Traditionally this component records the distance the car has traveled since the trip meter was set. This functionality is usually a result of the simple systems in place, and is tied to the revolutions of the wheels on the vehicle. As this system has more hardware at it's disposal, the trip meter can implement more functionality than a standard trip meter, making it much more useful in examining the performance of the car. Figure \ref{trip} shows the trip meter panel being displayed on the screen.

\figuremacro{trip}{Trip Meter Panel}{This figure shows two independent trip meters and the best record speed data}

The first unique point of this trip meter, is that it contains two independent meters. This is useful as it allows the driver to evaluate the statistics of two overlapping trips. One trip meter can be used to record the distance traveled since the car was last charged, while the other can be used to record the distance traveled in the last week or month. This independence allows the operator to decide how best to use the trip meter data, resulting in a high level of flexibility. In figure \ref{trip} the trip meters are located on the left, sitting above each other.

Each trip meter records the distance traveled, the time elapsed since the meter was started, the time the car has been moving since the meter was started and calculations based on the elapsed and moving time. These statistics are displayed live to the user, but are not logged, as the logging functionality  is taken car of by a different component in the system. The trip meter panel also displays the current moving speed in the top right. Below the moving speed are the best run records in seconds. This allows the driver to have quick feedback as to how the car is performing, without having to do lots of processing on logged data.

\subsection{Distance Driven}
\label{sec:distancedriven}
The most important part of a trip meter, is the distance that the meter has recorded. This is shown on figure \ref{trip} at the bottom of each trip meter. The meter stores the distance driven internally as a double length floating point number, but displays it on the screen as a rounded integer. This is done in order improve precision for later calculations as other values will depend on the distance that has been driven. Equation \ref{eq:distancedriven} shows the the distance is calculated based on the GPS position.

\begin{align}
\label{eq:distancedriven}
\mathrm{Distance} &=\mathrm{Distance}_\mathrm{lastrun} + \mathrm{Speed}(\mathrm{Time}_\mathrm{now}-\mathrm{Time}_\mathrm{lastrun})
\end{align}

The method for working out is a continuous function that is based on the last known distance the car has traveled. This method was chosen as it does not require any information other than the last time the formula was run, and the last distance calculated. This also makes the trip meter flexible in that the distance calculation does not need to be processed at exact intervals. If messages are dropped for whatever reason, the calculation will still take place, though it will not be as accurate as it could be. The calculation will be able to cope with fluctuations in the message timing, and can adapt to the speed of the GPS being used.

The downside of this method, is that big changes in time can cause problems with the calculation. If the signal drops out for an extended period of time, such as going through a tunnel, the calculation in \ref{eq:distancedriven} would have a big margin for error. In order to prevent this, the trip meter will ignore large time differences. If two calculations are over 10 seconds apart, the result will not be trusted, and not be used in the calculation. This prevents GPS signal loss from having an adverse effect on the trip meter calculations, but does impose a limit on the trip meter.

The limitation of this method of calculating the distance driven is that it relies on the GPS messages being sent to it. If the GPS signal is lost, the distance driven will not be increased. This is a limitation imposed by the use of the GPS sensor, and cannot be avoided, as attempting to guess the distance driven while the GPS signal has been lost has a very high probability of being incorrect.

\subsection{Time Elapsed}

Another variable displayed on the trip panel is the time elapsed. This is simply the time elapsed since the trip meter was last reset. This time increments even while the GPS signal has been lost, performing a stop-watch like action on the trip. While it may seem natural to just record the time that the trip meter was started and subtract it from the current system time, this would lead to problems during the system startup. The time needs to increment even while the GPS is connecting, and must be resistant to changes in the systems internal clock. Thus the time is calculated similar to section \ref{sec:distancedriven}. The formula used to calculate the elapsed time is given in equation \ref{eq:timeelapsed}. The time elapsed is displayed as the highest element of the trip meter in Figure \ref{trip}

\begin{align}
\label{eq:timeelapsed}
\mathrm{TimeElapsed} &=\mathrm{TimeElapsed}_\mathrm{lastrun} + (\mathrm{Time}_\mathrm{now}-\mathrm{Time}_\mathrm{lastrun})
\end{align}

Much like the distance, this method of calculation depends on the last known values. This means it does not matter at the actual time the system trip started once the timer has been running. This makes it resistant to changes in the system time, and thus makes the timer more robust. 
This timer records the time elapsed on the nanosecond level, as it is used elsewhere in calculations. For display, the timer converts these values into the traditional hours, minutes, seconds format that is easy for the operator to read.

\subsection{ Moving Time}

An aspect of the cars telemetry that would be interesting to the driver is the cars moving time. This is defined as the time in which the car has spent in motion. The main use of this data point is to contrast it against the elapsed time, to highlight how long the car has spent sitting still in traffic. This variable is also useful to record for future calculations, such as working out the average speed of the trip. This element is calculated according to equation \ref{eq:timeelapsed}, except that it will not update if the current speed of the car is 0~km/h. As such this element requires the speed of the car to be processed, so it cannot be calculated when the GPS signal is lost. This fits in with the functionality defined in section \ref{sec:distancedriven}, as the trip meter will not update the moving time or distance driven if the GPS signal is lost. The moving time is displayed below the elasped time and above the distance driven in Figure \ref{trip}

\subsection{Average Speed}

When reviewing the trip meter data, it is useful to know the average speed the car was traveling during the trip. Having this information allows the driver to better understand the characteristics of the drive. This element is also easy to calculate, as the time elapsed and the distance driven are already available. Equation \ref{eq:averagespeed} shows the formula used to calculate the average speed. The calculated value is displayed to the right of the elapsed time in figure \ref{trip}.

\begin{align}
\label{eq:averagespeed}
\mathrm{Average Speed} &=\frac{\mathrm{Distance Driven}}{\mathrm{Time Elapsed}}
\end{align}

This value is calculated whenever the distance driven or elapsed time values are updated. As this value is using two calculated values, it does not need to worry about discrepancies in time or the loss of the GPS signal.

\subsection{Average Moving Speed}

The average moving speed is like the average speed. The only difference is that it uses the moving time to calculate the speed, rather than the elapsed time. This is done using the same equation \ref{eq:averagespeed}, only substituting "Time Elapsed" for "Moving Time". This value provides the average speed of the car when it was actually being driven, thus ignoring time spent waiting in traffic.

\subsection{Reset}

The final functionality of each independent trip meter is the reset button. This button resets the trip meter it is attached to. The distance driven, moving and elapsed time counters will all display zero, and the average speed calculators will display zero. As each trip meter is independent, one can be reset without affecting the other. To reset the trip meters, the driver just has to press the reset button, located below the average moving time and to the right of the distance driven in Figure \ref{trip}.

\subsection{Current Speed}

The trip meters provide statistics on where the car has been driven, but do not provide much insight into the instantaneous speed of the vehicle. As this variable is already being used in calculations it is trivial to display it to the driver. This is displayed using a simple digit element, and appears in the top right corner of Figure \ref{trip}.


% ---------------------------------------------------------------------------
%: ----------------------- end of thesis sub-document ------------------------
% ---------------------------------------------------------------------------



%: ----------------------- contents from here ------------------------

\section{Interial Measurement Unit Display Panel}

Developed in section 


\figuremacro{imu}{IMU display panel}{This figure shows the information output by the IMU daemon }



% ---------------------------------------------------------------------------
%: ----------------------- end of thesis sub-document ------------------------
% ---------------------------------------------------------------------------



%: ----------------------- contents from here ------------------------

\section{Trip Meter}



% ---------------------------------------------------------------------------
%: ----------------------- end of thesis sub-document ------------------------
% ---------------------------------------------------------------------------



%: ----------------------- contents from here ------------------------

\section{Trip Meter}



% ---------------------------------------------------------------------------
%: ----------------------- end of thesis sub-document ------------------------
% ---------------------------------------------------------------------------



%: ----------------------- contents from here ------------------------

\section{About}

% ---------------------------------------------------------------------------
%: ----------------------- end of thesis sub-document ------------------------
% ---------------------------------------------------------------------------




%: ----------------------- contents from here ------------------------

\section{Settings}

The previous features have all discussed elements of the interface that are useful to any general driver. These are features that are useful while operating the vehicle. There also exists another user that has a different set of requirements to fulfill. This user is the maintainer of the car. This user will need access to various aspects for debugging and development purposes, including the ability to easily examine and transfer the internal logs of the system. In order to help simplify the navigation of the program, these elements are located in the settings screen. A screen-shot of the settings screen is shown in figure \ref{options}.

\figuremacro{options}{Options Panel}{This figure shows the options panel, which displays various utilities that are needed by the maintainer of the vehicle.}


% ---------------------------------------------------------------------------
%: ----------------------- end of thesis sub-document ------------------------
% ---------------------------------------------------------------------------



%: ----------------------- contents from here ------------------------

\section{Debug}

The top left most button in figure \ref{options} will navigate to the debug panel. This panel is shown in figure \ref{messages}. The purpose of this panel is to display the messages that are being received from the daemon programs, see section \ref{sec:design}. This panel displays the byte level contents of each message to the user as it is received. This is used to help facilitate debugging, as the writer of the daemon is able to see exactly what is being received in the interface.

\figuremacro{messages}{Debug Messages Panel}{This figure shows the Debug panel, it allows the user to filter which messages are displayed and alternates the colour between each received message}

This display panel has various buttons that exist on the top of the screen. These buttons allow known messages to be filtered out, enabling specific debugging to be undertaken. As discussed earlier, the first 3 bytes of a message are used to filter different kinds of messages. When a button is enabled, it will allow those messages to be drawn on the screen. Multiple combinations of messages can be combined to display any combination. The buttons labeled, TBS,GPS,IMU and DIO will filter the ASCII values of the first three bytes of each message. The all button will ignore any filtering, and will allow any message to be displayed, including ones that may not have been developed or seen before.

This panel uses a ConsoleElement to display it's data as seen in section \ref{sec:consoleelement}. This allows new messages to be displayed on the bottom of the screen, and travel upwards as they become outdated. It also provides a clear distinction between messages, as seen in figure \ref{messages}. The messages alternate between red and white as they are displayed, making it easy to identify the trailing portions of each message.

\section{Network Status Display}

In order to verify that the system is working correctly, it is helpful to view the messages being sent as in the Debug Panel. This panel can provide too much information, and requires the user to understand hex in order to verify that each message is being received. In order to simplify this, a panel was developed that displays the status of each different type of message. This uses the simple boolean display element previously mentioned in section \ref{sec:diopanel}. This element will display green whenever messages are currently being received by the system. If no message is received in the previous 10 seconds, the marker will transition to red in order to illustrate this. This allows a quick view as to the health of the system, and highlights components that need to have further inspection.

\figuremacro{netstat}{Daemon Status Panel}{This figure shows the status of the various daemons, green indicates the daemon is active, red indicates a timeout}

\section{System Logs Display}

As mentioned in section \ref{sec:dmesglogging}, all aspects of this system report errors by logging them to the syslog application. This prints out the warnings and errors created to \emph{/var/log/messages}. This panel loads and displays any information logged into this file. This information may come from any of the daemons mentioned in section \ref{sec:design}, the interface itself, or even other systems that are running on the Linux installation.

\figuremacro{dmesg}{Syslog Panel}{This figure shows the System Log panel, this displays the most recent entries in \emph{/var/log/messages}}

An important aspect of this panel is that it only displays the most recent information that has been logged. This makes it compatible with the ConsoleElement described in section \ref{sec:consoleelement}. The Console Element will automatically take care of displaying the last pieces of information it has received. This would produce the desired functionality, however further optimization needed to be undertaken.

The file \emph{/var/log/messages}, is logged to by all applications since the system started. Thus it grows in size the longer the system has been operational. Without optimizing the way this file is read, it would take O(n) time to display the newest entries in the file on the Console Element, where n is the number of entries in the file. This is not acceptable as the time taken to display the logs will increase the longer the application is running. In order to fix this problem, the file is read from the end of the file minus a 2048 byte offset.  The maximum width that can be displayed is 60 characters, and the maximum possible lines that can be shown on the screen is 20. This yields the result of 1200 possible characters as seen in equation \ref{eq:maximumdisplaycharacters}. This size is roughly doubled to 2048, in order to ensure more than enough lines are read for each update. This method of reading the last 2048 characters and finding the unique lines inside them is not affected by the file size, so the time taken to update the display is now O(1). The update is not dependent on the number of lines in \emph{/var/log/messages}

\begin{align}
\label{eq:maximumdisplaycharacters}
\mathrm{Maximum\ Characters} &= \mathrm{Characters\ per\ line}* \mathrm{Number\ of \ Lines} \\
&= 20*60 \notag \\
&=1200 \notag 
\end{align}

As the actual log file may be updated very quickly the screen must not update automatically. This is to stop the screen transitioning while the user may be reading something on it. As such the panel does not support any automatically updating functionality. In order to update the screen, a system like that in section \ref{sec:screenabout} is used. A invisible button is created to span the whole display. When pressed, this button will read the \emph{/var/log/messages} file and update the screen accordingly. This allows for the maximum utilization of space possible, as no space needs to be created to display a refresh button. The messages will also be updated when this panel becomes the active panel displayed on the screen.


% ---------------------------------------------------------------------------
%: ----------------------- end of thesis sub-document ------------------------
% ---------------------------------------------------------------------------



%: ----------------------- contents from here ------------------------

\section{Copy Log Files}

Discussed in section \ref{sec:filelogger} is a utility that records the telemetry data locally. In order to make the files created by this logger easier to access, the user interface provides a mechanism for them to be copied onto an external usb drive. Figure \ref{copy} shows the display asking for user confirmation to copy the log files. If the user confirms the action, the program will navigate to the directory in which the log files are stored, and copy them byte by byte to the external drive. Once completed the screen will display the figure \ref{copycompleted}. 
\figuremacro{copy}{Copy Popup}{Popup asking the user if they wish to copy log files to the an external drive}

If an error occurs, the system will log the results to the in-built system logger, syslog. This will allow the results to be viewable at \emph{/var/log/messages}, and thus can be viewed by the Syslog panel. While copying, the interface will display the progress of the action. It does this by indicating the amount of files to be copied, and the current file it is copying across. This is implemented so the user can be sure the system has not hung or encountered another error.

\figuremacro{copycompleted}{Copy Completed Popup}{This popup confirms that all the log files have been correctly copied to the usb drive}

\section{Delete Log Files}

In order to prevent the storage media from filling up, and to help organize the files when they are copied over the interface provides a mechanism to delete all the log files. Figure \ref{delete} shows a popup asking the user if they wish to delete the files. This functionality functions much the same way as copy, except that it does not require an external drive. It will simply open the folder and delete the files one by one. This action will also display its progress, and will provide the user with a message once the action is complete.

\figuremacro{delete}{Delete Log Files Popup}{This asks the user if they want to delete all the log files stored on the system}

% ---------------------------------------------------------------------------
%: ----------------------- end of thesis sub-document ------------------------
% ---------------------------------------------------------------------------


\section{Exit Program}

In order to allow other programs that use the various hardware, such as the screen or the touch driver, to run the current application must be exited. The system facilitates this be placing an exit action on the options panel. When this button is pressed, the popup shown in figure \ref{exit} will be displayed. This popup confirms that the user wishes to exit the application. If the user agrees, the program will exit and the screen will be set to black. If the user cancels, the system will return to the previously displayed panel, allowing normal operation to resume.

\figuremacro{exit}{Exit Popup}{This popup confirms that the user wants to exit the running application}




% ---------------------------------------------------------------------------
%: ----------------------- end of thesis sub-document ------------------------
% ---------------------------------------------------------------------------



% this file is called up by thesis.tex
% content in this file will be fed into the main document

\chapter{Performance} % top level followed by section, subsection


% ----------------------- paths to graphics ------------------------

% change according to folder and file names
\ifpdf
    \graphicspath{{7/figures/PNG/}{7/figures/PDF/}{7/figures/}}
\else
    \graphicspath{{7/figures/EPS/}{7/figures/}}
\fi


% ----------------------- contents from here ------------------------


\section{Limitations}

In testing the final system, issues arose which limited the depth in which the systems could be evaluated. As such only partial evaluation has been undertaken. This was done in an effort to provide some indication as to the stability of the system, and this section will provide analysis over what data could be obtained.

\subsection{Restricted Access to Vehicle}
\label{sec:access}
On the 15th of August the vehicle that was the main target for development was involved in a car accident. For safety reasons, the author was not allowed to work on the car in any form after this accident. This severely limited the amount of testing that could take place, as many aspects of the hardware are installed in the car which was no longer accessible.

\section{Message Timings}

This section reveals the time in between messages for the various components that transmit them. This was done using an eyebot to run the programs, and the messages were received on another laptop. Errors in the measurement can be contributed to the process switching active on both devices, as well as the fact that the laptop was communicating via a wireless network connection. All tests were conducted for approximately 20 minutes.

\subsection{Arduino Board}

Figure \ref{arduinotiming} shows the time between messages received vs the running time in the program. It is interesting to note that there are occasions when the time between two messages can be as much as 4 seconds. As the graph displays, this phenomenon appears to occur every 7 seconds. This is most likely caused by delays in the operating system on the receiving or sending end, as the process will not switch to the foreground automatically. As the delay is rarely more than 4 seconds, this will not interfere with the timeouts implemented in section \ref{toolkitfilter}

\figuremacro{arduinotiming}{Inter message timing for Arduino network messages}{The time delay between messages plotted against the time the program has been running for}

\subsection{Accelerometer}

The accelerometer runs at a much faster rate than the Arduino board. This is reflected in graph \ref{imutiming}. It can easily be seen that this program sends out messages much more frequently.  It still suffers the same problem with longer delays occurring. However each occurrence is still well under any timeout limits. 

\figuremacro{imutiming}{Inter message timing for accelerometer network messages}{The time delay between messages plotted against the time the program has been running for}

\subsection{GPS and Battery Monitor}

Due to both being physically installed in the vehicle, it was not possible to evaluate the message timings of these systems. See section \ref{sec:access}


\section{Memory Utilization}

This section shows plots of the memory usage of each daemon vs the running time of the program. If a program has a memory leak, it will eventually cause exhaust the systems memory and cause instability. These graphs exist to highlight any issues that may occur with running programs.

\subsection{Arduino Board}

Sadly, as shown in figure \ref{memoryarduino}, the Arduino based daemon does have a minor memory leak. After running the program for one hour, the memory usage had increased by 100kb. This memory leak is not too much of a problem. The system will never be on for more than a day at one time. After a day running non-stop, the program will only consume a little over 2mb. This is nothing compared to the 64mb of ram available to the system, so this problem is very minor.


\figuremacro{memoryarduino}{Memory usage of the Arduino daemon}{Memory usage (in kb) plotted against the time the program has been running for}

\subsection{Accelerometer}

As shown in figure \ref{memoryimu} the Accelerometer daemon has a memory leak at the rate of 760kb per hour. This is much worse than the leak in the Arduino board. Extrapolating this usage out to one day yields a memory usage of approximately 18mb. This is still small enough to not cause major problems, though it does warrant an investigation into the cause of the memory leak. As long as the system is not run for over a day without a restart, this leak will not be a problem.

\figuremacroW{memoryimu}{Memory usage of the Accelerometer daemon}{Memory usage (in kb) plotted against the time the program has been running for}{0.8}

\subsection{User interface}

The user interface was tested with a simulator sending messages from all four components developed in section \ref{sec:design}. As figure \ref{memorygui} shows, the user interface does not have any memory leaks. The interface can run for any amount of time and will not cause system instability. This is a good result as the user interface will always remain responsive.

\figuremacroW{memorygui}{Memory usage of the user interface}{Memory usage (in kb) plotted against the time the program has been running for}{0.8}

\subsection{GPS and Battery Monitor}

Due to both being physically installed in the vehicle, it was not possible to evaluate the message timings of these systems. See section \ref{sec:access}. Running these programs without access to the hardware would not yield reliable results.




% ---------------------------------------------------------------------------
% ----------------------- end of thesis sub-document ------------------------
% ---------------------------------------------------------------------------               % discussion of results


% this file is called up by thesis.tex
% content in this file will be fed into the main document

\chapter{Conclusions} % top level followed by section, subsection


% ----------------------- paths to graphics ------------------------

% change according to folder and file names
\ifpdf
    \graphicspath{{8/figures/PNG/}{8/figures/PDF/}{8/figures/}}
\else
    \graphicspath{{8/figures/EPS/}{8/figures/}}
\fi

% ----------------------- contents from here ------------------------



This project shows that it is possible to use high level programming techniques when developing for embedded devices. The hardware used in this project was powerful enough to separate into multiple components. This is significant as it provides flexibility in the development of embedded devices and highlights how these products are becoming more like traditional computers.

This project also illustrates the viability of using existing networking techniques, such as TCP/IP, in providing communication between different devices. The application of this technology to embedded systems provides increased robustness and decreases development time.  This shows that a network of small embedded systems that each contribute a portion of the overall system functionality is a viable solution to the problems of scaling and system cost.

The ability to automatically log all the information recorded by the car is of great help to the goals of the REV team. The group will be able to log much more detailed data, and the system has been constructed in a way to make it easily extend-able. This will ultimately help in the research and development of the vehicle the system is running on.

\section{Limitations}

The system currently has a few limitations. These are caused by various constraints in the hardware and the design of the system.

The system will be unable to support too many more additions with it's current processing power. The CPU is not able to run many high frequency data sources at the same time. This will limit the future expansion unless new hardware is acquired.

The system also suffers some issues with the memory utilization of some components. Further investigation will need to take place in order to eliminate the bugs that are causing these issues.

Due to time constraints the Arduino and Accelerometer were not developed fully in regards to the user interface. This is a limitation that can be overcome in the future, when a more specific use of these components is required.

\section{Advantages}

The system is easily expendable, providing a framework for the transmission and display of any data that the group may want to develop in the future. The development of the windowing toolkit makes it much easier to work with the touch interface on the device. 

The messaging protocol developed, and the use of a message abstraction, makes it much easier to communicate between different sections. The fact that the complicated components of networking are abstracted from the developer make it easier to conceptualize what is happening. A component in this system only needs to worry about the data being sent, and not a complicated method of ensuring it is received correctly.

Due to it's distributed nature, the system is easy and cost-effective to extend. Rather than having to source expensive new hardware to perform all the tasks, cheaper hardware can be used to take the role of some of the tasks. This allows the complete system to adapt to the conditions that it is being used in.

\section{Applications}

The system has applications in the project it was developed for. It provides a clean user interface for the display of data, and logging functionality so it can be analyzed later.

The system is a proof of concept that embedded systems can communicate reliably using network protocols. This concept could be applied to newly developed systems, enabling products that can be purchased in sections, with each section providing different or improved functionality. This paves the way for embedded systems to be configured via pre-built components , rather than developed using one monolithic structure. 

\section{Future work}

The first aspect of future work is the finishing of the Arduino and Accelerometer components. These components are able to be decoded, but do not have very useful user interfaces. 

Currently the networking is not encrypted, which does expose the possibility that it can be attacked. This does require access to the network, which is not easily, and understanding of the protocols being communicated. Future work would include the investigation of SSL in order to prevent any third party from viewing the data that is transmitted.

The windowing toolkit could be further developed by providing rendering of the full alphabet at varying sizes. This would simplify the use of the toolkit, as it would be much easier to make text appear. The toolkit could also be further developed into it's own standalone package, either with dependencies on the existing libraries or not.

Navigation could be added to the maps portion. A path-finding algorithm could be developed and implemented that would be able to direct the driver towards a specific destination.

Further research could be conducted into sensor fusion. As the system has access to GPS, the speedometer pulse frequency, current, as well as the acceleration the car is undertaking a system could be developed in order to compensate for loss of a sensor. For instance the other sensors could approximate the GPS position when inside a tunnel. This would increase the user experience, has it would appear that the GPS works in all locations.

New hardware should also be investigated for the system. Either in the form of additional devices to offset the load, or replacement of the main component running the system. This would increase the processing power of the system, and allow more complicated functionality to be developed.


% ---------------------------------------------------------------------------
%: ----------------------- end of thesis sub-document ------------------------
% ---------------------------------------------------------------------------



 






        % description of lab methods




% --------------------------------------------------------------
%:                  BACK MATTER: appendices, refs,..
% --------------------------------------------------------------

% the back matter: appendix and references close the thesis


%: ----------------------- bibliography ------------------------

% The section below defines how references are listed and formatted
% The default below is 2 columns, small font, complete author names.
% Entries are also linked back to the page number in the text and to external URL if provided in the BibTex file.

% PhDbiblio-url2 = names small caps, title bold & hyperlinked, link to page 
\begin{multicols}{2} % \begin{multicols}{ # columns}[ header text][ space]
\begin{tiny} % tiny(5) < scriptsize(7) < footnotesize(8) < small (9)

\bibliographystyle{Latex/Classes/PhDbiblio-url2} % Title is link if provided
\renewcommand{\bibname}{References} % changes the header; default: Bibliography

\bibliography{9_backmatter/beautrepp} % adjust this to fit your BibTex file

\end{tiny}
\end{multicols}

% --------------------------------------------------------------
% Various bibliography styles exit. Replace above style as desired.

% in-text refs: (1) (1; 2)
% ref list: alphabetical; author(s) in small caps; initials last name; page(s)
%\bibliographystyle{Latex/Classes/PhDbiblio-case} % title forced lower case
%\bibliographystyle{Latex/Classes/PhDbiblio-bold} % title as in bibtex but bold
%\bibliographystyle{Latex/Classes/PhDbiblio-url} % bold + www link if provided

%\bibliographystyle{Latex/Classes/jmb} % calls style file jmb.bst
% in-text refs: author (year) without brackets
% ref list: alphabetical; author(s) in normal font; last name, initials; page(s)

%\bibliographystyle{plainnat} % calls style file plainnat.bst
% in-text refs: author (year) without brackets
% (this works with package natbib)


% --------------------------------------------------------------

% according to Dresden med fac summary has to be at the end
%
% Thesis Abstract -----------------------------------------------------


%\begin{abstractslong}    %uncommenting this line, gives a different abstract heading
\begin{abstracts}        %this creates the heading for the abstract page

The topic of electric vehicles is becoming increasingly popular due to rising fuel costs and growing concern over emissions. Despite this attention, most electric vehicles have little or no telemetry systems, making many aspects of there operation and efficiency a mystery. 

The aim of this project was to develop an extend-able system in order to capture various data-points that can be available in a vehicle, as well as an interface to display this data inside the car. The design developed differs from traditional embedded systems by being completely modular. It uses existing network protocols to allow the system to be distributed between various smaller embedded components. This will enable it to be easily extended, should the need for more data-points arise, and allows the use of many smaller systems to be implemented incrementally, rather than one expensive monolithic design.

In order to help facilitate robustness and code re-use, a windowing toolkit was also developed. This provides a common platform for user interaction, and defines pre-built components in order to speed up the design and implementations of the user interface.

By exposing and recording more data, deeper analysis can be done on the efficiency of the car, and help justify different technological improvements to the vehicle. The higher granularity of data acquired can also be used to analyze the economy of the vehicle in different conditions.

\end{abstracts}
%\end{abstractlongs}


% ---------------------------------------------------------------------- 


%: Declaration of originality

% Thesis statement of originality -------------------------------------

% Depending on the regulations of your faculty you may need a declaration like the one below. This specific one is from the medical faculty of the university of Dresden.

\begin{declaration}        %this creates the heading for the declaration page

I herewith declare that I have produced this paper without the prohibited assistance of third parties and without making use of aids other than those specified; notions taken over directly or indirectly from other sources have been identified as such. This paper has not previously been presented in identical or similar form to any other German or foreign examination board.

The thesis work was conducted from XXX to YYY under the supervision of PI at ZZZ.

\vspace{10mm}

CITY,


\end{declaration}


% ----------------------------------------------------------------------

% this file is called up by thesis.tex
% content in this file will be fed into the main document

%: ----------------------- name of chapter  -------------------------
\appendix

 % top level followed by section, subsection


%: ----------------------- paths to graphics ------------------------

% change according to folder and file names
\ifpdf
    \graphicspath{{9_backmatter/figures/PNG/}{9_backmatter/figures/PDF/}{9_backmatter/figures/}}
\else
    \graphicspath{{9_backmatter/figures/EPS/}{9_backmatter/figures/}}
\fi

%: ----------------------- contents from here ------------------------

\chapter{Remotely Accessing Eyebot}

\section{SSH Details}

To connect to the eyebot, issue the command from any linux machine that has ssh installed.

\begin{verbatim}
ssh root@192.168.1.50
\end{verbatim}

You will then be prompted for a password. There is no password set by default. Hit enter when prompted for a password. The table \ref{app:tab:sshdetails} shows the settings to connect.

\begin{table}
    \begin{tabular}{|l|l|}
        \hline
        Property        & Value        \\ \hline
        SSH Address(IP) & 192.168.1.50 \\ 
        Account name    & root         \\ 
        Password        & <empty>      \\
        \hline
    \end{tabular}
	\caption{SSH Details}
	\label{app:tab:sshdetails}
\end{table}

\section{Installing programs to Eyebot}

Each program is accompanied with a makefile. This makefile will connect to the default IP address of the machine, and install the program in the folder \emph{/root/demo/rev2011/}. To install any program just run 

\begin{verbatim}
make install
\end{verbatim}

From the source folder of the desired program on the development machine. Press enter when prompted for a password.

\chapter{Development environment setup}

\section{Buildroot setup}

The following section is based upon the legacy buildroot installation guide obtainable from

\begin{verbatim}
http://docwiki.gumstix.org/Buildroot
\end{verbatim}

This section purpose is to describe the steps needed to be undertaken to build the compiler and libraries required by the eyebot. The online guide requires some tweaks for newer versions of GCC to compile the older versions. This section will detail the changes required for Ubuntu 10.10. Other distributions may require further changes aswell.

\subsection{Checkout Buildroot source code}

\begin{enumerate}
	\item To perform this section you must have subversion installed on your machine. This is a common version control system, and it's installation is out of the scope of this document.
	\item Checkout the 1083 release of buildroot. The eyebot was built with older buildroot libraries, these are version 1083 in the repository. To check out this source use the command (all one line)

\begin{verbatim}
	svn co -r1083 http://svn.gumstix.com/gumstix-buildroot/trunk 
	gumstix-buildroot
\end{verbatim}

\subsection{Setup environment}

	\item Next you will have to setup the buildroot make environment. Change into the folder buildroot was downloaded too.
	\item Remove the existing config (if present) with

\begin{verbatim}
	rm ./.config
\end{verbatim}

	\item Run the autoconfig program

\begin{verbatim}
	make defconfig
\end{verbatim}

	You will be prompted to select settings about the eyebot. The eyebot contains a 400mhz xscale processor, so select the options according to this. Any other options can be selected with the defaults.

\subsection{Modify the source}
	\item Modifying the source files.

	The source files will not compile correctly on ubuntu 10.10 without some minor modifications. This has occured because most systems use gcc4+, while this code was originally compiled on gcc3

\begin{enumerate}
	\item Makeinfo will complain about not being installed to fix this edit the file. 
\begin{verbatim}
	toolchain_build_arm_nofpu/binutils-2.17-build/Makefile 
\end{verbatim}
	Change line 258 to 
\begin{verbatim}
	MAKEINFO = “/usr/bin/makeinfo”
\end{verbatim}
	or whatever the correct path to makeinfo is.

	\item The older GCC did not require permission settings when opening a file, this must be added in. Edit the file
\begin{verbatim}
	toolchain_build_arm_nofpu/gcc-3.4.5/gcc/collect2.c
\end{verbatim}
	On line 1537 change it to read
\begin{verbatim}
	redir_handle = open (redir, O_WRONLY | O_TRUNC | O_CREAT,0666);
\end{verbatim}

	\item The limits.h file is not included by default when compiling using the newer GCC (caused by changes in libraries) to fix edit.
\begin{verbatim}
	build_arm_nofpu/linux-2.6.17gum/scripts/mod/sumversion.c
\end{verbatim}
	And insert at the top of the file
\begin{verbatim}
	#include <limits.h>
\end{verbatim}
\end{enumerate}

	\item Include the e2fs files for the UUID functions. UUID functions are required by the ZeroMQ library. Run
\begin{verbatim}
	make menuconfig
\end{verbatim}
	And navigate the menus and enable the e2fs library.

\subsection{Build}
	\item Compile the environment. Run 
\begin{verbatim}
	make 
\end{verbatim}
	This will download and compile the build environment. Note that some packages may no longer be mirrored online. You will need to search the internet and manually download these files to the dl folder that is created. The compile process will fail with the name of the missing packages, so you can hunt them down.

\end{enumerate}

\section{Compiling ZeroMQ for the arm cpu}

This section details the commands used to build zeroMQ libraries to run on the arm board. These commands work with zeroMQ version 2.2 on a ubuntu 10.10 machine. 

\begin{enumerate}
	\subsection{Setup path variable}
	\item First you must setup the path variable so the automake process knows the location of the compiler. To do this run this command in the terminal. (all one line)
 \begin{verbatim}
	export PATH=$PATH:%PATHTOBUILDROOT%/gumstix-buildroot/
	build_arm_nofpu/staging_dir/bin/
\end{verbatim}
	Where \%PATHTOBUILDROOT% is the path to the buildroot folder

	\subsection{Run Automake}
	\item ZeroMQ contains an configure utility that will generate make files for you. Run this utility with this command (all one line)
\begin{verbatim}
CC=arm-linux-uclibc-gcc CFLAGS=" -pthread -D_REENTRANT
 -march=armv5te -mtune=xscale -DZMQ_FORCE_POLL" 
CXX=arm-linux-uclibc-g++ 
CXXFLAGS=" -pthread -D_REENTRANT -march=armv5te -mtune=xscale
 -DZMQ_FORCE_POLL"
LDFLAGS="-L%X%/gumstix-buildroot/build_arm_nofpu/e2fsprogs-1.38/lib/"
CPPFLAGS="-I%X%/gumstix-buildroot/build_arm_nofpu/staging_dir/include 
-I%X%/gumstix-buildroot/build_arm_nofpu/e2fsprogs-1.38/lib/uuid/"
./configure --host=arm-linux 
--prefix=%X%/build_arm_nofpu/root
\end{verbatim}
	Where \%X\% is the path to the gumstix buildroot environment setup earlier.

	\subsection{Build}
	\item To build ZeroMQ run the command
\begin{verbatim}
	make
\end{verbatim}
	This will build zeroMQ. To obtain the libraries run
\begin{verbatim}
	make installl
\end{verbatim}
	This will copy the libraries to \emph{\%X\%/build\_arm\_nofpu/root/\%ZEROMQVER\%/}
\end{enumerate}

\chapter{Performance Testing Scripts}

\section{Memory Utilization}

In section \ref{sec:memutilization} the memory usage of various programs was recorded. To do this the following script was used

\begin{verbatim}
while true
do
	date +%s | awk '{ printf "%s ", $0 }' >> output.txt;
	ps | grep %EXECUTABLE_NAME% | grep -v grep | head -n1 \
	| awk '{print $3}' | awk '{printf "%s", $0}' >> output.txt;
	echo '' >> output.txt
	sleep 20;
done

\end{verbatim}

\chapter{Utility Programs}
\label{app:sec:utilities}

During the development of this system, various utility programs were developed in order to help debug and test the system.

\section{ZeroMQ Viewer}

This component is located in the file 

\begin{verbatim}
	Utilities/programs/ZeroMQViewer
\end{verbatim}

This program is a subscriber that will attach to any publisher component. It print the time-stamp and hexadecimal values of the message whenever a message is received.

\section{ZeroMQ Emulator}

This component is located at 

\begin{verbatim}
	Utilities/programs/ZeroMQMessageEmulator/
\end{verbatim}

This component acts as a publisher, it will transmit messages that can be interpreted by subscribers. It loads the information ASCII information from a text file and interprets it as hexadecimal. It supports spaces in-between hexadecimal values and comments so the file should be human readable. Examples are contained in the folder with the binary.

\section{ZeroMQ Evaluator}

This program is contained in 

\begin{verbatim}
	Utilities/programs/ZeroMQMessageEvaluator/
\end{verbatim}

This program is a modified version of the Viewer. It will read in the specified number of messages and then stop. It also records the start and end time that this took place. This allows the time it take a certain volume or messages to be transmitted to be recorded and analyzed.

\chapter{Message Timing Tables}
\label{app:sec:timing}

The following section shows the data measured for the time taken to receive varying numbers of messages. This data was obtained by running the Evaluate program on the eyebot, and sending messages from a laptop. This laptop generated messages as fast as it could using the Emulator program. When the Evaluate program recieved the desired number of messages it would halt and print the timing results.  The results are available in system clock ticks, and the system time  that elapsed during this period. The results are also available in varying message sizes. These messages were generated randomly using \emph{/dev/random} on a linux system, and stored in a file that the Emulator could load and run.

\begin{table}
    \begin{tabular}{*{7}{|l}|}
        \hline
Test Number & 1 & 2 & 3 & 4 & 5 &  ~ \\  \hline \hline
1000 & 20000 & 40000 & 20000 & 20000 & 20000 & ~\\ 
10000 & 170000 & 180000 & 120000 & 160000 & 160000 & ~ \\
100000 & 870000 & 970000 & 1250000 & 1010000 & 1000000 & ~  \\ 
1000000 & 9560000 & 10220000 & 10230000 & 11220000 & 9710000 & ~ \\ \hline
Test Number &  6 & 7 & 8 & 9 & 10 &  Average \\  \hline \hline
1000  & 10000 & 20000 & 30000 & 10000 & 30000 & 22000 \\ 
10000  & 140000 & 160000 & 130000 & 120000 & 150000 & 149000 \\ 
100000 &  1020000 & 1040000 & 940000 & 890000 & 1050000 & 1004000 \\ 
1000000 & 9580000 & 10980000 & 10570000 & 10800000 & 11760000 & 10463000 \\
        \hline
    \end{tabular}
	\caption{Clock ticks Vs Number of messages for 32byte random data}
\end{table}

\begin{table}
    \begin{tabular}{*{7}{|l}|}
        \hline
Test Number & 1 & 2 & 3 & 4 & 5 &  ~ \\  \hline \hline
1000  & 243750 & 212332 & 256312 & 3156460 & 341224 & ~\\
10000 & 1419011 & 1349485 & 1436313 & 1534302 & 1349040 & ~ \\
100000 & 15145656 & 15590986 & 15703572 & 15657899 & 15649278 & ~ \\
1000000 & 152007642 & 156742950 & 157926991 & 158742238 & 157040941 & ~ \\ \hline
Test Number &  6 & 7 & 8 & 9 & 10 &  Average \\  \hline \hline
1000 & 356874 & 204477 & 386403 & 233102 & 359293 & 575022.7 \\
10000 & 1405605 & 1662092 & 1697169 & 1501454 & 1398061 & 1475253.2 \\
100000& 15663356 & 15267476 & 15416116 & 15363871 & 15265640 & 15472385 \\
1000000 & 156889964 & 158464172 & 153960703 & 159458560 & 158363673 & 156959783.4 \\
	\hline
    \end{tabular}
\caption{Time (us) Vs Number of messages for 32byte random data}
\end{table}


\begin{table}
    \begin{tabular}{*{7}{|l}|}
        \hline
Test Number & 1 & 2 & 3 & 4 & 5 &  ~ \\  \hline \hline
10 & 30000 & 20000 & 20000 & 20000 & 0 & ~ \\
100 & 70000 & 30000 & 90000 & 60000 & 50000 & ~ \\
1000 & 80000 & 30000 & 70000 & 40000 & 80000 & ~ \\
10000 & 150000 & 190000 & 160000 & 210000 & 150000 & ~ \\
100000 & 1930000 & 2210000 & 1850000 & 1890000 & 2080000 & ~ \\ \hline
Test Number &  6 & 7 & 8 & 9 & 10 &  Average \\  \hline \hline
10 & 20000 & 30000 & 60000 & 40000 & 30000 & 27000 \\
100 & 10000 & 40000 & 10000 & 30000 & 10000 & 40000 \\
1000 & 50000 & 30000 & 50000 & 50000 & 80000 & 56000 \\
10000& 280000 & 290000 & 180000 & 220000 & 180000 & 201000 \\
100000 & 2150000 & 2030000 & 1870000 & 2020000 & 2090000 & 2012000 \\
        \hline
    \end{tabular}
	\caption{Clock ticks Vs Number of messages for 128byte random data}
\end{table}

\begin{table}
    \begin{tabular}{*{7}{|l}|}
        \hline
Test Number & 1 & 2 & 3 & 4 & 5 &  ~ \\  \hline \hline
10 & 323135 & 306990 & 220649 & 220649 & 191386 & ~ \\
100 & 393113 & 205562 & 315971 & 328754 & 357496 & ~ \\
1000 & 496815 & 438725 & 389029 & 505927 & 465334 & ~ \\
10000 & 2740074 & 2723108 & 2601828 & 2576363 & 2711420 & ~ \\
100000 & 27669065 & 27824854 & 26609896 & 25919529 & 26034126 & ~ \\ \hline
Test Number &  6 & 7 & 8 & 9 & 10 &  Average \\  \hline \hline
10 & 266575 & 220406 & 317913 & 307281 & 189971 & 256495.5 \\
100 & 301369 & 291968 & 198676 & 311395 & 331863 & 303616.7 \\
1000 & 488639 & 5468 & 519439 & 474776 & 523848 & 430800 \\
10000 & 2732068 & 2710042 & 2631575 & 2639119 & 3595522 & 2766111.9 \\
100000 & 27232333 & 26982161 & 27233195 & 26655991 & 27697161 & 26985831.1 \\
	\hline
    \end{tabular}
\caption{Time (us) Vs Number of messages for 128byte random data}
\end{table}



\begin{table}
    \begin{tabular}{*{7}{|l}|}
        \hline
Test Number & 1 & 2 & 3 & 4 & 5 &  ~ \\  \hline \hline
10 & 0 & 20000 & 0 & 20000 & 30000 & ~ \\
100 & 10000 & 50000 & 50000 & 30000 & 10000 & ~ \\
1000 & 90000 & 70000 & 50000 & 50000 & 80000 & ~ \\
10000 & 230000 & 170000 & 190000 & 270000 & 350000 & ~ \\
100000 & 2170000 & 1800000 & 2230000 & 2240000 & 2020000 & ~ \\ \hline
Test Number &  6 & 7 & 8 & 9 & 10 &  Average \\  \hline \hline
 10 & 30000 & 20000 & 20000 & 10000 & 40000 & 19000 \\
100 & 20000 & 40000 & 90000 & 30000 & 30000 & 36000 \\
1000 & 10000 & 30000 & 10000 & 40000 & 70000 & 50000 \\ 
10000 & 260000 & 260000 & 150000 & 170000 & 3177625 & 522762.5 \\
100000 & 2180000 & 2040000 & 1980000 & 1940000 & 2260000 & 2086000 \\
        \hline
    \end{tabular}
	\caption{Clock ticks Vs Number of messages for 512byte random data}
\end{table}

\begin{table}
    \begin{tabular}{*{7}{|l}|}
        \hline
Test Number & 1 & 2 & 3 & 4 & 5 &  ~ \\  \hline \hline
10 & 202744 & 299140 & 243188 & 224886 & 310320 & ~ \\
100 & 3134384 & 336945 & 344085 & 229151 & 226242 & ~ \\
1000 & 527391 & 546638 & 584975 & 529211 & 548475 & ~ \\
10000 & 3049670 & 2849869 & 2671973 & 2945053 & 3007291 & ~ \\
100000 & 30712493 & 31676926 & 30686329 & 30285899 & 31565685 & ~ \\ \hline
Test Number &  6 & 7 & 8 & 9 & 10 &  Average \\  \hline \hline
10 & 320544 & 306020 & 220204 & 372283 & 325050 & 282437.9 \\
100 & 203972 & 377633 & 368549 & 319270 & 336170 & 587640.1 \\ 
1000 & 477733 & 595697 & 484988 & 587090 & 561919 & 544411.7 \\
10000 & 2872847 & 2891389 & 2926986 & 3177625 & 2871762 & 2926446.5 \\
100000 & 30283039 & 31059473 & 31532265 & 32491691 & 31771619 & 31206541.9 \\
	\hline
    \end{tabular}
\caption{Time (us) Vs Number of messages for 512byte random data}
\end{table}



\begin{table}
    \begin{tabular}{*{7}{|l}|}
        \hline
Test Number & 1 & 2 & 3 & 4 & 5 &  ~ \\  \hline \hline
10 & 30000 & 20000 & 30000 & 20000 & 30000 & ~ \\
100 & 100000 & 40000 & 40000 & 40000 & 0 & ~ \\
1000 & 60000 & 30000 & 60000 & 50000 & 90000 & ~ \\
10000 & 140000 & 200000 & 190000 & 200000 & 260000 & ~ \\
100000 & 1800000 & 1890000 & 2410000 & 2210000 & 1700000 & ~ \\ \hline
Test Number &  6 & 7 & 8 & 9 & 10 &  Average \\  \hline \hline
10 & 10000 & 10000 & 30000 & 10000 & 30000 & 22000 \\
100 & 80000 & 30000 & 10000 & 60000 & 20000 & 42000 \\
1000 & 60000 & 50000 & 60000 & 0 & 80000 & 54000 \\
10000 & 330000 & 200000 & 170000 & 280000 & 210000 & 218000 \\
100000 & 2210000 & 2180000 & 2410000 & 2060000 & 1900000 & 2077000 \\
        \hline
    \end{tabular}
	\caption{Clock ticks Vs Number of messages for 1024byte random data}
\end{table}

\begin{table}
    \begin{tabular}{*{7}{|l}|}
        \hline
Test Number & 1 & 2 & 3 & 4 & 5 &  ~ \\  \hline \hline
10 & 344441 & 307364 & 330483 & 315630 & 344759  & ~ \\
100 & 383476 & 351497 & 350538 & 417289 & 339274  & ~ \\
1000 & 539349 & 601257 & 716008 & 548840 & 711210  & ~ \\
10000 & 3429580 & 3497672 & 3247387 & 4596539 & 3276066  & ~ \\
100000 & 36594732 & 35013510 & 34448621 & 34661021 & 35992978  & ~ \\ \hline
Test Number &  6 & 7 & 8 & 9 & 10 &  Average \\  \hline \hline
10 & 310805 & 231331 & 267536 & 345342 & 335819 & 313351 \\
100 & 352329 & 366358 & 197007 & 312781 & 215691 & 328624 \\
1000 & 566045 & 520301 & 500614 & 591720 & 548868 & 584421.2 \\
10000 & 3489945 & 3114163 & 3573906 & 3328063 & 3611572 & 3516489.3 \\
100000 & 35032912 & 35081602 & 33674882 & 34417770 & 35152012 & 35007004 \\
	\hline
    \end{tabular}
\caption{Time (us) Vs Number of messages for 1024byte random data}
\end{table}


\begin{table}
    \begin{tabular}{*{7}{|l}|}
        \hline
Test Number & 1 & 2 & 3 & 4 & 5 &  ~ \\  \hline \hline
10 & 30000 & 30000 & 40000 & 40000 & 20000 & ~ \\
100 & 80000 & 40000 & 70000 & 30000 & 70000 & ~ \\
1000 & 80000 & 50000 & 20000 & 60000 & 60000 & ~ \\
10000 & 170000 & 200000 & 230000 & 190000 & 180000 & ~ \\
100000 & 1730000 & 1580000 & 1580000 & 1570000 & 1490000 & ~ \\ \hline
Test Number &  6 & 7 & 8 & 9 & 10 &  Average \\  \hline \hline
10 & 40000 & 10000 & 70000 & 30000 & 40000 & 35000 \\
100 & 50000 & 60000 & 30000 & 20000 & 80000 & 53000 \\
1000 & 20000 & 60000 & 20000 & 60000 & 30000 & 46000 \\
10000 & 210000 & 230000 & 190000 & 170000 & 220000 & 199000 \\
100000 & 1970000 & 1570000 & 2100000 & 2070000 & 1900000 & 1756000 \\
        \hline
    \end{tabular}
	\caption{Clock ticks Vs Number of messages for 10kilobyte random data}
\end{table}

\begin{table}
    \begin{tabular}{*{7}{|l}|}
        \hline
Test Number & 1 & 2 & 3 & 4 & 5 &  ~ \\  \hline \hline
10 & 377608 & 331011 & 285319 & 186330 & 224346  & ~ \\
100 & 387303 & 395092 & 349759 & 344359 & 364674  & ~ \\
1000 & 606900 & 694347 & 651131 & 470944 & 584394  & ~ \\
10000 & 3034906 & 3020777 & 2916406 & 3016008 & 3014637  & ~ \\
100000 & 29623959 & 30323297 & 31451455 & 31070588 & 30427494  & ~ \\ \hline
Test Number &  6 & 7 & 8 & 9 & 10 &  Average \\  \hline \hline
10 & 383401 & 243644 & 345962 & 264199 & 319124 & 296094.4 \\
100 & 337517 & 395891 & 387481 & 216318 & 401400 & 357979.4 \\
1000 & 491953 & 581445 & 603782 & 583448 & 704321 & 597266.5 \\
10000 & 2909019 & 3082147 & 3222074 & 3322947 & 2926152 & 3046507.3 \\
100000 & 29400026 & 31617928 & 30323380 & 30336067 & 32110036 & 30668423 \\
	\hline
    \end{tabular}
\caption{Time (us) Vs Number of messages for 10kilobyte random data}
\end{table}


\begin{table}
    \begin{tabular}{*{5}{|l}|}
        \hline

Message Size & 128 & 512 & 1024 & 10240 \\ \hline \hline
10 & 256495.5 & 282437.9 & 313351 & 213073.6 \\
100 & 303616.7 & 587640.1 & 328624 & 304995.2 \\
1000 & 430800 & 544411.7 & 584421.2 & 390158.225 \\
10000 & 2766111.9 & 2926446.5 & 3516489.3 & 2304761.925 \\
100000 & 26985831.1 & 31206541.9 & 35007004 & 23324844.25 \\
        \hline
    \end{tabular}
	\caption{Clock ticks Vs Number of messages combined data for varying message sizes}
\end{table}

\begin{table}
    \begin{tabular}{*{5}{|l}|}
        \hline
 Message Size & 128 bytes & 512 bytes & 1024 bytes & 10 kilobytes \\ \hline \hline
10 & 256.4955 & 282.4379 & 313.351 & 213.0736 \\
100 & 303.6167 & 587.6401 & 328.624 & 304.9952 \\
1000 & 430.8 & 544.4117 & 584.4212 & 390.158225 \\
10000 & 2766.1119 & 2926.4465 & 3516.4893 & 2304.761925 \\
100000 & 26985.8311 & 31206.5419 & 35007.004 & 23324.84425 \\
	\hline
    \end{tabular}
\caption{Time(ms) Vs Number of messages combined data for varying message sizes}
\end{table}


% ---------------------------------------------------------------------------
%: ----------------------- end of thesis sub-document ------------------------
% ---------------------------------------------------------------------------



\end{document}
